% Options for packages loaded elsewhere
\PassOptionsToPackage{unicode}{hyperref}
\PassOptionsToPackage{hyphens}{url}
%
\documentclass[
]{article}
\usepackage{amsmath,amssymb}
\usepackage{iftex}
\ifPDFTeX
  \usepackage[T1]{fontenc}
  \usepackage[utf8]{inputenc}
  \usepackage{textcomp} % provide euro and other symbols
\else % if luatex or xetex
  \usepackage{unicode-math} % this also loads fontspec
  \defaultfontfeatures{Scale=MatchLowercase}
  \defaultfontfeatures[\rmfamily]{Ligatures=TeX,Scale=1}
\fi
\usepackage{lmodern}
\ifPDFTeX\else
  % xetex/luatex font selection
\fi
% Use upquote if available, for straight quotes in verbatim environments
\IfFileExists{upquote.sty}{\usepackage{upquote}}{}
\IfFileExists{microtype.sty}{% use microtype if available
  \usepackage[]{microtype}
  \UseMicrotypeSet[protrusion]{basicmath} % disable protrusion for tt fonts
}{}
\makeatletter
\@ifundefined{KOMAClassName}{% if non-KOMA class
  \IfFileExists{parskip.sty}{%
    \usepackage{parskip}
  }{% else
    \setlength{\parindent}{0pt}
    \setlength{\parskip}{6pt plus 2pt minus 1pt}}
}{% if KOMA class
  \KOMAoptions{parskip=half}}
\makeatother
\usepackage{xcolor}
\usepackage[margin=1in]{geometry}
\usepackage{graphicx}
\makeatletter
\def\maxwidth{\ifdim\Gin@nat@width>\linewidth\linewidth\else\Gin@nat@width\fi}
\def\maxheight{\ifdim\Gin@nat@height>\textheight\textheight\else\Gin@nat@height\fi}
\makeatother
% Scale images if necessary, so that they will not overflow the page
% margins by default, and it is still possible to overwrite the defaults
% using explicit options in \includegraphics[width, height, ...]{}
\setkeys{Gin}{width=\maxwidth,height=\maxheight,keepaspectratio}
% Set default figure placement to htbp
\makeatletter
\def\fps@figure{htbp}
\makeatother
\setlength{\emergencystretch}{3em} % prevent overfull lines
\providecommand{\tightlist}{%
  \setlength{\itemsep}{0pt}\setlength{\parskip}{0pt}}
\setcounter{secnumdepth}{-\maxdimen} % remove section numbering
% definitions for citeproc citations
\NewDocumentCommand\citeproctext{}{}
\NewDocumentCommand\citeproc{mm}{%
  \begingroup\def\citeproctext{#2}\cite{#1}\endgroup}
\makeatletter
 % allow citations to break across lines
 \let\@cite@ofmt\@firstofone
 % avoid brackets around text for \cite:
 \def\@biblabel#1{}
 \def\@cite#1#2{{#1\if@tempswa , #2\fi}}
\makeatother
\newlength{\cslhangindent}
\setlength{\cslhangindent}{1.5em}
\newlength{\csllabelwidth}
\setlength{\csllabelwidth}{3em}
\newenvironment{CSLReferences}[2] % #1 hanging-indent, #2 entry-spacing
 {\begin{list}{}{%
  \setlength{\itemindent}{0pt}
  \setlength{\leftmargin}{0pt}
  \setlength{\parsep}{0pt}
  % turn on hanging indent if param 1 is 1
  \ifodd #1
   \setlength{\leftmargin}{\cslhangindent}
   \setlength{\itemindent}{-1\cslhangindent}
  \fi
  % set entry spacing
  \setlength{\itemsep}{#2\baselineskip}}}
 {\end{list}}
\usepackage{calc}
\newcommand{\CSLBlock}[1]{\hfill\break\parbox[t]{\linewidth}{\strut\ignorespaces#1\strut}}
\newcommand{\CSLLeftMargin}[1]{\parbox[t]{\csllabelwidth}{\strut#1\strut}}
\newcommand{\CSLRightInline}[1]{\parbox[t]{\linewidth - \csllabelwidth}{\strut#1\strut}}
\newcommand{\CSLIndent}[1]{\hspace{\cslhangindent}#1}
\ifLuaTeX
\usepackage[bidi=basic]{babel}
\else
\usepackage[bidi=default]{babel}
\fi
\babelprovide[main,import]{spanish}
% get rid of language-specific shorthands (see #6817):
\let\LanguageShortHands\languageshorthands
\def\languageshorthands#1{}
\usepackage{amsfonts}
\usepackage{booktabs}
\usepackage{longtable}
\usepackage{array}
\usepackage{float}
\usepackage{pdflscape}
\usepackage{threeparttable}
\usepackage[normalem]{ulem}
\usepackage{amssymb}
\usepackage{amsmath}
\usepackage{geometry}
\usepackage{booktabs}
\usepackage{longtable}
\usepackage{array}
\usepackage{multirow}
\usepackage{wrapfig}
\usepackage{float}
\usepackage{colortbl}
\usepackage{pdflscape}
\usepackage{tabu}
\usepackage{threeparttable}
\usepackage{threeparttablex}
\usepackage[normalem]{ulem}
\usepackage{makecell}
\usepackage{xcolor}
\ifLuaTeX
  \usepackage{selnolig}  % disable illegal ligatures
\fi
\usepackage{bookmark}
\IfFileExists{xurl.sty}{\usepackage{xurl}}{} % add URL line breaks if available
\urlstyle{same}
\hypersetup{
  pdftitle={Matemáticas 11° - Segundo Periodo 2023},
  pdflang={es},
  hidelinks,
  pdfcreator={LaTeX via pandoc}}

\title{Matemáticas 11° - Segundo Periodo 2023}
\usepackage{etoolbox}
\makeatletter
\providecommand{\subtitle}[1]{% add subtitle to \maketitle
  \apptocmd{\@title}{\par {\large #1 \par}}{}{}
}
\makeatother
\subtitle{Informe profesional de evaluación}
\author{true}
\date{15 de enero de 2025}

\begin{document}
\maketitle

{
\setcounter{tocdepth}{3}
\tableofcontents
}
\section*{Introducción}\label{introducciuxf3n}
\addcontentsline{toc}{section}{Introducción}

\section{Pregunta 01}\label{pregunta-01}

\begin{center}\rule{0.5\linewidth}{0.5pt}\end{center}

\includegraphics{/home/manjaro_lenovo/Documentos/Math/2023-Matematicas-11-2/01/01.png}

\includegraphics{/home/manjaro_lenovo/Documentos/Math/2023-Matematicas-11-2/01/all2.png}

\begin{center}\rule{0.5\linewidth}{0.5pt}\end{center}

\subsection{Análisis}\label{anuxe1lisis}

La encuesta realizada a 120 personas revela que:

\begin{itemize}
\tightlist
\item
  30\% de las personas (36 personas) están interesadas en adoptar un
  loro.
\item
  40\% de las personas (48 personas) están interesadas en adoptar un
  perro.
\item
  30\% de las personas (36 personas) están interesadas en adoptar un
  gato.
\end{itemize}

Ahora, evaluamos las opciones proporcionadas:

\begin{itemize}
\tightlist
\item
  \textbf{Opción A:} Presenta una tabla con los números correctos (36
  para loros, 48 para perros y 36 para gatos).
\item
  \textbf{Opción B:} Presenta un gráfico de barras con porcentajes
  incorrectos, por lo que está descartada.
\item
  \textbf{Opción C:} Muestra una lista con los números correctos, aunque
  el formato es diferente.
\item
  \textbf{Opción D:} Presenta un gráfico de barras con los números
  correctos.
\end{itemize}

Aunque las opciones A, C y D tienen la información correcta, la opción A
es la más clara y directa, ya que presenta los datos en una tabla
sencilla y fácil de entender.

\subsubsection{\texorpdfstring{\textbf{Respuesta:} La opción correcta es
\textbf{A}.}{Respuesta: La opción correcta es A.}}\label{respuesta-la-opciuxf3n-correcta-es-a.}

\begin{center}\rule{0.5\linewidth}{0.5pt}\end{center}

\subsubsection{Análisis de la Pregunta 2023-Matematicas-11-2-01, según
SABER 11
ICFES}\label{anuxe1lisis-de-la-pregunta-2023-matematicas-11-2-01-seguxfan-saber-11-icfes}

\subsection{Características de la
Pregunta}\label{caracteruxedsticas-de-la-pregunta}

\begin{itemize}
\tightlist
\item
  Nivel de Desempeño: 2
\item
  Competencia: Interpretación y Representación

  \begin{itemize}
  \tightlist
  \item
    Afirmación: Comprende y transforma la información cuantitativa y
    esquemática presentada en distintos formatos
  \item
    Evidencia: Da cuenta de las características básicas de la
    información presentada en diferentes formatos como series, gráficas,
    tablas y esquemas
  \end{itemize}
\item
  Componente: Aleatorio
\item
  Estándar Asociado: Interpreto, produzco y comparo representaciones
  gráficas adecuadas para presentar diversos tipos de datos (diagramas
  de barras, diagramas circulares)
\item
  ¿Qué evalúa?: La capacidad de identificar la representación gráfica
  que corresponde correctamente a un conjunto de datos presentados en
  forma verbal
\end{itemize}

\subsection{Solución y
Justificación}\label{soluciuxf3n-y-justificaciuxf3n}

La información dada establece que de 120 personas encuestadas:

\begin{itemize}
\tightlist
\item
  \(30\% = 36\) personas adoptarían un loro
\item
  \(40\% = 48\) personas adoptarían un perro\\
\item
  \(30\% = 36\) personas adoptarían un gato
\end{itemize}

La opción A presenta una tabla que muestra exactamente estos valores
numéricos de forma clara y ordenada, siendo la representación más
precisa y directa de los datos.

Respuesta correcta: A

\subsection{Análisis de Distractores}\label{anuxe1lisis-de-distractores}

\begin{itemize}
\tightlist
\item
  Opción B: El diagrama de barras muestra porcentajes que no
  corresponden a los datos originales
\item
  Opción C: Aunque muestra los valores correctos en un diagrama
  circular, la distribución visual puede ser más difícil de interpretar
\item
  Opción D: El diagrama de barras horizontal, si bien tiene los valores
  correctos, presenta la información de forma menos directa que una
  tabla
\end{itemize}

\subsection{Contenidos Matemáticos
Curriculares}\label{contenidos-matemuxe1ticos-curriculares}

\begin{itemize}
\tightlist
\item
  Estadística
\item
  Período 1

  \begin{itemize}
  \tightlist
  \item
    Diferentes tipos de representación

    \begin{itemize}
    \tightlist
    \item
      Tablas univariadas
    \item
      Diagrama de barras individuales
    \item
      Diagrama circular
    \end{itemize}
  \end{itemize}
\end{itemize}

\subsection{Clasificación}\label{clasificaciuxf3n}

\begin{itemize}
\tightlist
\item
  Genérico: Sí
\item
  Eje Axial Disciplinar: Univariadas - Interpretación de tablas y barras
  elementales
\item
  Tarea: Identificar la representación gráfica que corresponde
  correctamente a un conjunto de datos presentados textualmente
\item
  Grado sugerido: 7°
\end{itemize}

\section{Analisis de Resultados}\label{analisis-de-resultados}

\subsection{Competencia 1: Pensamiento
Numérico}\label{competencia-1-pensamiento-numuxe9rico}

\section{Pregunta 02}\label{pregunta-02}

\begin{center}\rule{0.5\linewidth}{0.5pt}\end{center}

\includegraphics{/home/manjaro_lenovo/Documentos/Math/2023-Matematicas-11-2/02/01.png}

\begin{center}\rule{0.5\linewidth}{0.5pt}\end{center}

\subsection{Análisis}\label{anuxe1lisis-1}

Para determinar cuál de las expresiones es equivalente al polinomio
\(p(x)\) dado para valores de \(x\) distintos de 0, debemos simplificar
el polinomio original y compararlo con las opciones proporcionadas.

El polinomio original es:

\[ p(x) = \frac{\frac{1}{100} x^4 - \frac{12}{10} x^3 + 36x^2 + 100x}{10x} \]

Vamos a simplificar esta expresión paso a paso.

\begin{enumerate}
\def\labelenumi{\arabic{enumi}.}
\item
  \textbf{Simplificar el numerador y el denominador:}

  El numerador es:

  \[ \frac{1}{100} x^4 - \frac{12}{10} x^3 + 36x^2 + 100x \]

  El denominador es:

  \[ 10x \]
\item
  \textbf{Dividir cada término del numerador por el denominador:}

  \[
  p(x) = \frac{\frac{1}{100} x^4}{10x} - \frac{\frac{12}{10} x^3}{10x} + \frac{36x^2}{10x} + \frac{100x}{10x}
  \]
\item
  \textbf{Simplificar cada término:}

  \[
  \frac{\frac{1}{100} x^4}{10x} = \frac{1}{1000} x^3
  \]

  \[
  \frac{\frac{12}{10} x^3}{10x} = \frac{12}{100} x^2 = \frac{3}{25} x^2
  \]

  \[
  \frac{36x^2}{10x} = \frac{36}{10} x = 3.6x
  \]

  \[
  \frac{100x}{10x} = 10
  \]
\item
  \textbf{Combinar los términos simplificados:}

  \[
  p(x) = \frac{1}{1000} x^3 - \frac{3}{25} x^2 + 3.6x + 10
  \]
\item
  \textbf{Convertir las fracciones a decimales para facilitar la
  comparación:}

  \[
  \frac{1}{1000} = 0.001
  \]

  \[
  \frac{3}{25} = 0.12
  \]

  Por lo tanto, la expresión simplificada es:

  \[
  p(x) = 0.001x^3 - 0.12x^2 + 3.6x + 10
  \]
\item
  \textbf{Comparar con las opciones dadas:}

  \begin{itemize}
  \tightlist
  \item
    \textbf{A.}
    \(\frac{1}{1.000} x^3 - \frac{12}{100} x^2 + \frac{36}{10} x + 10\)
  \item
    \textbf{B.}
    \(\frac{1}{1.000} x^4 - \frac{12}{100} x^3 + \frac{36}{10} x^2 + 10x\)
  \item
    \textbf{C.} \(\frac{1}{100} x^3 - \frac{12}{10} x^2 + 36x + 100\)
  \item
    \textbf{D.}
    \(\frac{10}{100} x^4 - \frac{120}{10} x^3 + \frac{36}{10} x^2 + 10x\)
  \end{itemize}

  La opción correcta es la \textbf{A}, ya que coincide con la expresión
  simplificada:

  \[
  p(x) = \frac{1}{1.000} x^3 - \frac{12}{100} x^2 + \frac{36}{10} x + 10
  \]
\end{enumerate}

\subsubsection{Conclusión}\label{conclusiuxf3n}

La opción correcta es la \textbf{A}.

\begin{center}\rule{0.5\linewidth}{0.5pt}\end{center}

\subsubsection{Análisis de la Pregunta 2023-Matematicas-11-2-02, según
SABER 11
ICFES}\label{anuxe1lisis-de-la-pregunta-2023-matematicas-11-2-02-seguxfan-saber-11-icfes}

\begin{itemize}
\tightlist
\item
  Nivel de Desempeño: 4
\item
  Competencia: Interpretación y Representación

  \begin{itemize}
  \tightlist
  \item
    Afirmación: Comprende y transforma la información cuantitativa y
    esquemática presentada en distintos formatos
  \item
    Evidencia: Transforma la representación de una o más piezas de
    información
  \end{itemize}
\item
  Componente: Numérico-Variacional
\item
  Estándar Asociado: Utilizo las técnicas de aproximación en procesos
  infinitos numéricos.
\item
  ¿Qué evalúa?: La capacidad para simplificar expresiones algebraicas
  racionales y determinar expresiones equivalentes.
\item
  Respuesta Correcta: A
\item
  Justificación: Para encontrar la expresión equivalente, debemos
  simplificar \(p(x)\):
\end{itemize}

\begin{enumerate}
\def\labelenumi{\arabic{enumi})}
\item
  El polinomio original es:
  \(p(x) = \frac{\frac{1}{100}x^4 - \frac{12}{10}x^3 + 36x^2 + 100x}{10x}\)
\item
  Simplificando término a término:
\end{enumerate}

\begin{itemize}
\tightlist
\item
  \(\frac{\frac{1}{100}x^4}{10x} = \frac{1}{1000}x^3\)
\item
  \(-\frac{\frac{12}{10}x^3}{10x} = -\frac{12}{100}x^2\)
\item
  \(\frac{36x^2}{10x} = \frac{36}{10}x\)
\item
  \(\frac{100x}{10x} = 10\)
\end{itemize}

\begin{enumerate}
\def\labelenumi{\arabic{enumi})}
\setcounter{enumi}{2}
\tightlist
\item
  Por lo tanto:
  \(p(x) = \frac{1}{1000}x^3 - \frac{12}{100}x^2 + \frac{36}{10}x + 10\)
\end{enumerate}

Esta expresión coincide con la opción A.

\begin{itemize}
\tightlist
\item
  Distractores:

  \begin{itemize}
  \tightlist
  \item
    B: No realiza la división por 10x correctamente
  \item
    C: Usa denominadores incorrectos
  \item
    D: Multiplica en lugar de dividir por 10x
  \end{itemize}
\item
  Contenidos Matemáticos Curriculares:

  \begin{itemize}
  \tightlist
  \item
    Álgebra y Cálculo
  \item
    Período: 2

    \begin{itemize}
    \tightlist
    \item
      Álgebra

      \begin{itemize}
      \tightlist
      \item
        Expresiones algebraicas
      \item
        Operaciones con expresiones algebraicas
      \end{itemize}
    \end{itemize}
  \end{itemize}
\item
  No genérico
\item
  Eje Axial Disciplinar: Eje 3 (Términos Semejantes, Factorización
  Básica: Álgebra)
\item
  Tarea: Simplificar una expresión algebraica racional para determinar
  una expresión equivalente
\item
  Grado: 11°
\end{itemize}

\subsection{Competencia 2: Pensamiento
Geométrico}\label{competencia-2-pensamiento-geomuxe9trico}

\section{Pregunta 03}\label{pregunta-03}

\begin{center}\rule{0.5\linewidth}{0.5pt}\end{center}

\includegraphics{/home/manjaro_lenovo/Documentos/Math/2023-Matematicas-11-2/03/01.png}

\begin{center}\rule{0.5\linewidth}{0.5pt}\end{center}

\subsection{Análisis}\label{anuxe1lisis-2}

Paula dibujó un triángulo PQR, y se proporciona la siguiente información
sobre las razones trigonométricas de sus ángulos:

\begin{itemize}
\tightlist
\item
  \(\sin(P) > 0\), \(\sin(Q) > 0\), \(\sin(R) > 0\)
\item
  \(\cos(P) > 0\), \(\cos(Q) > 0\), \(\cos(R) < 0\)
\end{itemize}

\begin{enumerate}
\def\labelenumi{\arabic{enumi}.}
\item
  \textbf{Seno positivo:}\\
  El seno es positivo para ángulos entre 0° y 180°, lo cual es siempre
  cierto en un triángulo.
\item
  \textbf{Coseno positivo y negativo:}

  \begin{itemize}
  \tightlist
  \item
    \(\cos(P) > 0\) y \(\cos(Q) > 0\) indican que ambos ángulos \(P\) y
    \(Q\) son agudos (menos de 90°).
  \item
    \(\cos(R) < 0\) implica que el ángulo \(R\) es obtuso (más de 90°).
  \end{itemize}
\end{enumerate}

\subsubsection{Conclusión:}\label{conclusiuxf3n-1}

Dado que la suma de los ángulos en un triángulo es 180°, si \(R\) es
obtuso, entonces \(P\) y \(Q\) deben ser agudos. Por lo tanto, el
triángulo tiene un ángulo obtuso y se clasifica como
\textbf{obtusángulo}.

La opción correcta es:

\textbf{D. Obtusángulo; es decir, tiene un ángulo que mide más de 90°.}

\begin{center}\rule{0.5\linewidth}{0.5pt}\end{center}

\subsubsection{Análisis de la Pregunta 2023-Matematicas-11-2-03, según
SABER 11
ICFES}\label{anuxe1lisis-de-la-pregunta-2023-matematicas-11-2-03-seguxfan-saber-11-icfes}

\section{Análisis Pregunta
Matemáticas}\label{anuxe1lisis-pregunta-matemuxe1ticas}

\begin{itemize}
\tightlist
\item
  Nivel de Desempeño: 3
\item
  Competencia: Argumentación

  \begin{itemize}
  \tightlist
  \item
    Aprendizaje: Valida procedimientos y estrategias matemáticas
    utilizadas para dar solución a problemas.
  \item
    Evidencia: Plantea afirmaciones que sustentan o refutan una
    interpretación dada a la información disponible en el marco de la
    solución de un problema.
  \end{itemize}
\item
  Componente: Geométrico-Métrico
\item
  Estándar Asociado: Uso argumentos geométricos para resolver y formular
  problemas en contextos matemáticos y en otras ciencias.
\item
  ¿Qué evalúa?: La capacidad de clasificar un triángulo según sus
  ángulos usando información sobre razones trigonométricas.
\item
  Respuesta Correcta: D
\item
  Justificación: Dado que \(\cos(R) < 0\), el ángulo R debe ser obtuso
  (mayor a 90°), mientras que P y Q son agudos pues tienen coseno
  positivo. Por lo tanto, el triángulo es obtusángulo.
\item
  Distractores:

  \begin{itemize}
  \tightlist
  \item
    A: Confunde las condiciones del coseno negativo con ángulo recto
  \item
    B: No relaciona la información trigonométrica con los tipos de
    triángulos
  \item
    C: Interpreta erróneamente que cosenos positivos implican todos los
    ángulos agudos
  \end{itemize}
\item
  Contenidos Matemáticos Curriculares:

  \begin{itemize}
  \tightlist
  \item
    Geometría
  \item
    Período: 1

    \begin{itemize}
    \tightlist
    \item
      Teoría del ángulo (Clases, ángulos entre rectas paralelas)
    \end{itemize}
  \end{itemize}
\item
  No genérico
\item
  Eje Axial Disciplinar: Geometría del triángulo I
\item
  Tarea: Clasificar un triángulo según sus ángulos a partir de
  información sobre los signos de las razones trigonométricas de sus
  ángulos.
\item
  Grado: 10°
\end{itemize}

\subsection{Competencia 3: Pensamiento
Métrico}\label{competencia-3-pensamiento-muxe9trico}

\section{Pregunta 04}\label{pregunta-04}

\begin{center}\rule{0.5\linewidth}{0.5pt}\end{center}

\includegraphics{/home/manjaro_lenovo/Documentos/Math/2023-Matematicas-11-2/04/01.png}

\begin{center}\rule{0.5\linewidth}{0.5pt}\end{center}

\subsection{Análisis}\label{anuxe1lisis-3}

La tabla muestra la duración de las baterías de tres modelos de
celulares (X, Y y Z) en cinco pruebas diferentes. La tarea es comparar
las medianas de las duraciones de las baterías y determinar cuál de las
afirmaciones es correcta.

\subsubsection{Paso 1: Ordenar los datos para cada
modelo}\label{paso-1-ordenar-los-datos-para-cada-modelo}

\begin{itemize}
\item
  \textbf{Modelo X:} 570, 620, 550, 600, 560\\
  Ordenados: 550, 560, 570, 600, 620\\
  Mediana: 570 minutos
\item
  \textbf{Modelo Y:} 520, 690, 540, 670, 680\\
  Ordenados: 520, 540, 670, 680, 690\\
  Mediana: 670 minutos
\item
  \textbf{Modelo Z:} 520, 550, 620, 570, 710\\
  Ordenados: 520, 550, 570, 620, 710\\
  Mediana: 570 minutos
\end{itemize}

\subsubsection{Paso 2: Comparar las
medianas}\label{paso-2-comparar-las-medianas}

\begin{itemize}
\tightlist
\item
  Mediana de X: 570 minutos
\item
  Mediana de Y: 670 minutos
\item
  Mediana de Z: 570 minutos
\end{itemize}

\subsubsection{Paso 3: Evaluar las
opciones}\label{paso-3-evaluar-las-opciones}

\begin{itemize}
\item
  \textbf{Opción A:} Incorrecta. Z (570) no tiene la misma mediana que Y
  (670), y es menor que X (570), pero X y Z tienen la misma mediana.
\item
  \textbf{Opción B:} Incorrecta. Z (570) no tiene una mediana mayor que
  Y (670) ni que X (570).
\item
  \textbf{Opción C:} Correcta. Z (570) tiene la misma mediana que X
  (570) y ambas son menores que Y (670).
\item
  \textbf{Opción D:} Incorrecta. Z (570) no tiene una mediana mayor que
  X (570), ambas son iguales.
\end{itemize}

\subsubsection{Conclusión}\label{conclusiuxf3n-2}

La afirmación correcta es la \textbf{Opción C}.

\subsubsection{Respuesta final}\label{respuesta-final}

La opción correcta es \textbf{C.} El modelo Z tiene una mediana igual a
la del modelo X y menor que la del modelo Y.

\begin{center}\rule{0.5\linewidth}{0.5pt}\end{center}

\subsubsection{Análisis de la Pregunta 2023-Matematicas-11-2-04, según
SABER 11
ICFES}\label{anuxe1lisis-de-la-pregunta-2023-matematicas-11-2-04-seguxfan-saber-11-icfes}

\subsection{Nivel de Desempeño}\label{nivel-de-desempeuxf1o}

2

\subsection{Competencia}\label{competencia}

Interpretación y Representación

\begin{itemize}
\tightlist
\item
  Afirmación: Comprende y transforma la información cuantitativa y
  esquemática presentada en distintos formatos.
\item
  Evidencia: Da cuenta de las características básicas de la información
  presentada en diferentes formatos como series, gráficas, tablas y
  esquemas.
\end{itemize}

\subsection{Componente}\label{componente}

Aleatorio

\subsection{Estándar Asociado}\label{estuxe1ndar-asociado}

Interpreto nociones básicas relacionadas con el manejo de información
como población, muestra, variable aleatoria, distribución de
frecuencias, parámetros y estadígrafos.

\subsection{¿Qué evalúa?}\label{quuxe9-evaluxfaa}

La capacidad del estudiante para calcular e interpretar medidas de
tendencia central (mediana) a partir de datos presentados en una tabla y
realizar comparaciones entre estos valores.

\subsection{Respuesta Correcta y
Justificación}\label{respuesta-correcta-y-justificaciuxf3n}

La respuesta correcta es C.

Para cada modelo, ordenando los datos de menor a mayor:

Modelo X: \(550, 560, \mathbf{570}, 600, 620\) → Mediana = \(570\)

Modelo Y: \(520, 540, \mathbf{670}, 680, 690\) → Mediana = \(670\)

Modelo Z: \(520, 550, \mathbf{570}, 620, 710\) → Mediana = \(570\)

Por lo tanto:

\begin{itemize}
\tightlist
\item
  Mediana Z = Mediana X = \(570\)
\item
  Mediana Y = \(670\) \textgreater{} Mediana Z
\end{itemize}

\subsection{Motivos de las Opciones
Incorrectas}\label{motivos-de-las-opciones-incorrectas}

\begin{itemize}
\tightlist
\item
  A: Confunde las relaciones, pues Z no tiene mediana igual a Y
\item
  B: Invierte las relaciones de orden entre las medianas
\item
  D: Confunde la relación entre Z y X, que son iguales
\end{itemize}

\subsection{Contenidos Matemáticos
Curriculares}\label{contenidos-matemuxe1ticos-curriculares-1}

\begin{itemize}
\tightlist
\item
  Estadística

  \begin{itemize}
  \tightlist
  \item
    Período 2

    \begin{itemize}
    \tightlist
    \item
      Medidas de tendencia central

      \begin{itemize}
      \tightlist
      \item
        Mediana
      \end{itemize}
    \end{itemize}
  \end{itemize}
\end{itemize}

\subsection{¿Genérico o No
genérico?}\label{genuxe9rico-o-no-genuxe9rico}

Genérico

\subsection{Eje Axial Disciplinar}\label{eje-axial-disciplinar}

Eje 4: Descriptivos - Tendencia Central

\subsection{Tarea}\label{tarea}

Calcular y comparar las medianas de tres conjuntos de datos presentados
en una tabla.

\subsection{Grado}\label{grado}

10° - 11°

\subsection{Competencia 4: Pensamiento
Aleatorio}\label{competencia-4-pensamiento-aleatorio}

\section{Pregunta 05}\label{pregunta-05}

\begin{center}\rule{0.5\linewidth}{0.5pt}\end{center}

\includegraphics{/home/manjaro_lenovo/Documentos/Math/2023-Matematicas-11-2/05/all.png}

\begin{center}\rule{0.5\linewidth}{0.5pt}\end{center}

\subsection{Análisis}\label{anuxe1lisis-4}

Dado un segmento PQ, el objetivo es construir un nuevo segmento que sea
16 veces más largo que PQ, siguiendo un procedimiento similar al que se
usa para construir un segmento 4 veces más largo.

\subsubsection{Paso a paso:}\label{paso-a-paso}

\begin{enumerate}
\def\labelenumi{\arabic{enumi}.}
\item
  \textbf{Entendiendo el procedimiento inicial}: El proceso para
  construir un segmento 4 veces más largo que PQ probablemente involucra
  el uso de círculos para duplicar sucesivamente la longitud. Por
  ejemplo, un círculo permite duplicar la longitud, y repetir este
  proceso con otro círculo duplica nuevamente el segmento, alcanzando
  así una longitud 4 veces mayor.
\item
  \textbf{Escalando a 16 veces la longitud}: La longitud 16PQ es
  equivalente a realizar cuatro veces el proceso de duplicación (ya que
  16 = 2\^{}4). Si cada círculo permite duplicar la longitud, entonces
  serán necesarios cuatro círculos para llegar de PQ a 16PQ.
\item
  \textbf{Verificación de eficiencia}: Se consideraron otras
  posibilidades, como el uso de círculos para construir triángulos
  semejantes o métodos más avanzados. Sin embargo, no se encontró una
  forma más eficiente que reduzca el número de círculos necesario.
\end{enumerate}

\subsubsection{Conclusión:}\label{conclusiuxf3n-3}

La cantidad mínima de círculos necesaria para construir un segmento 16
veces más largo que PQ es \textbf{cuatro} (C).

\begin{center}\rule{0.5\linewidth}{0.5pt}\end{center}

\subsubsection{Análisis de la Pregunta 2023-Matematicas-11-2-05, según
SABER 11
ICFES}\label{anuxe1lisis-de-la-pregunta-2023-matematicas-11-2-05-seguxfan-saber-11-icfes}

\begin{itemize}
\tightlist
\item
  Nivel de Desempeño: 3
\item
  Competencia: Argumentación

  \begin{itemize}
  \tightlist
  \item
    Aprendizaje: Valida procedimientos y estrategias matemáticas
    utilizadas para dar solución a problemas
  \item
    Evidencia: Argumenta a favor o en contra de un procedimiento para
    resolver un problema a la luz de criterios presentados o
    establecidos
  \end{itemize}
\item
  Componente: Geométrico-Métrico
\item
  Estándar Asociado: Construyo y descompongo figuras y sólidos a partir
  de condiciones dadas
\item
  ¿Qué evalúa?: La capacidad de analizar un proceso geométrico de
  construcción y extenderlo a un caso más complejo
\end{itemize}

\subsection{Solución}\label{soluciuxf3n}

La respuesta correcta es C. 4 círculos

\subsubsection{Justificación}\label{justificaciuxf3n}

\begin{enumerate}
\def\labelenumi{\arabic{enumi}.}
\tightlist
\item
  En el procedimiento mostrado, cada par de círculos duplica la longitud
  del segmento inicial:

  \begin{itemize}
  \tightlist
  \item
    Con 2 círculos: \(PQ \rightarrow 2PQ\)
  \item
    Con 4 círculos: \(2PQ \rightarrow 4PQ\)
  \end{itemize}
\item
  Para llegar a 16 veces la longitud original, necesitamos duplicar 4
  veces:

  \begin{itemize}
  \tightlist
  \item
    \(2^4 = 16\)
  \item
    Cada duplicación requiere un círculo
  \item
    Por lo tanto: \(4\) círculos son necesarios
  \end{itemize}
\end{enumerate}

\subsubsection{Distractores}\label{distractores}

\begin{itemize}
\tightlist
\item
  A (1 círculo): Insuficiente, solo permitiría duplicar una vez
\item
  B (2 círculos): Solo alcanzaría para cuadruplicar el segmento
\item
  D (8 círculos): Excesivo, produciría un segmento de \(2^8 = 256\)
  veces el original
\end{itemize}

\subsection{Características}\label{caracteruxedsticas}

\begin{itemize}
\tightlist
\item
  Contenidos Matemáticos Curriculares:

  \begin{itemize}
  \tightlist
  \item
    Geometría
  \item
    Período 2

    \begin{itemize}
    \tightlist
    \item
      Polígonos y Transformaciones en el Plano

      \begin{itemize}
      \tightlist
      \item
        Homotecia
      \end{itemize}
    \end{itemize}
  \end{itemize}
\item
  No genérico
\item
  Eje Axial Disciplinar: Traslación, Rotación, Reflexión, Homotecias:
  Movimientos en el plano (x,y)
\item
  Tarea: Determinar la cantidad mínima de círculos necesarios para
  multiplicar por 16 la longitud de un segmento usando un método
  geométrico de duplicación sucesiva
\item
  Grado: 9°
\end{itemize}

\subsection{Competencia 5: Pensamiento
Variacional}\label{competencia-5-pensamiento-variacional}

\section{Pregunta 06}\label{pregunta-06}

\begin{center}\rule{0.5\linewidth}{0.5pt}\end{center}

\includegraphics{/home/manjaro_lenovo/Documentos/Math/2023-Matematicas-11-2/06/01.png}

\begin{center}\rule{0.5\linewidth}{0.5pt}\end{center}

\subsection{Análisis}\label{anuxe1lisis-5}

\subsection{Evaluación Sistemática}\label{evaluaciuxf3n-sistemuxe1tica}

\subsubsection{1. Análisis de la Ecuación
Original}\label{anuxe1lisis-de-la-ecuaciuxf3n-original}

La ecuación que deben resolver es: (x + 2)(x + 3) = 5(x + 3)

\subsubsection{2. Evaluación de Cada
Procedimiento}\label{evaluaciuxf3n-de-cada-procedimiento}

\paragraph{María}\label{maruxeda}

\begin{verbatim}
x² + 5x + 6 = 5x + 15       # Desarrolla el producto del lado izquierdo
x² + 6 = 15                 # Agrupa términos semejantes
x² = 15 - 6                 # Despeja x²
x² = 9                      # Simplifica
x = ±3                      # Obtiene las raíces
\end{verbatim}

\emph{Procedimiento correcto y completo}

\paragraph{Nelson}\label{nelson}

\begin{verbatim}
(x + 2)(x + 3) - 5(x + 3) = 0   # Pasa todo a un lado de la igualdad
(x + 3)[(x + 2) - 5] = 0        # Factoriza el término común (x + 3)
(x + 3)(x - 3) = 0              # Simplifica
x² - 9 = 0                      # Multiplica los términos
x = ±3                          # Obtiene las raíces
\end{verbatim}

\emph{Procedimiento correcto y completo}

\paragraph{Óscar}\label{uxf3scar}

\begin{verbatim}
x + 2 + x + 3 = 5 + x + 3       # Error: no desarrolla el producto
2x + 5 = x + 8                  # Agrupa términos
2x - x = 8 - 5                  # Despeja x
x = 3                           # Obtiene solo una raíz
\end{verbatim}

\emph{Procedimiento incorrecto}

\subsubsection{3. Conclusión}\label{conclusiuxf3n-4}

La respuesta correcta es la opción \textbf{B: Solo María y Nelson}.

Justificación: - María y Nelson obtienen correctamente ambas soluciones
(x = 3 y x = -3) - Sus procedimientos son algebraicamente válidos -
Óscar comete un error fundamental al no desarrollar el producto (x +
2)(x + 3) - Óscar solo obtiene una solución parcial y por un método
incorrecto

\subsection{Verificación}\label{verificaciuxf3n}

Para x = 3: - (3 + 2)(3 + 3) = 5(3 + 3) - (5)(6) = 5(6) - 30 = 30
\checkmark

Para x = -3: - (-3 + 2)(-3 + 3) = 5(-3 + 3) - (-1)(0) = 5(0) - 0 = 0
\checkmark

Las soluciones x = 3 y x = -3 son correctas, confirmando que María y
Nelson llegaron al resultado correcto por métodos válidos.

\begin{center}\rule{0.5\linewidth}{0.5pt}\end{center}

\subsubsection{Análisis de la Pregunta 2023-Matematicas-11-2-06, según
SABER 11
ICFES}\label{anuxe1lisis-de-la-pregunta-2023-matematicas-11-2-06-seguxfan-saber-11-icfes}

\subsubsection{Nivel de Desempeño: 3}\label{nivel-de-desempeuxf1o-3}

\subsubsection{Competencia:
Argumentación}\label{competencia-argumentaciuxf3n}

\begin{itemize}
\tightlist
\item
  Aprendizaje: Valida procedimientos y estrategias matemáticas
  utilizadas para dar solución a problemas.
\item
  Evidencia: Argumenta a favor o en contra de un procedimiento para
  resolver un problema a la luz de criterios presentados o establecidos.
\end{itemize}

\subsubsection{Componente:
Numérico-Variacional}\label{componente-numuxe9rico-variacional}

\subsubsection{Estándar Asociado:}\label{estuxe1ndar-asociado-1}

Justifico procedimientos algebraicos utilizados en la transformación de
expresiones algebraicas y verifico relaciones entre expresiones
algebraicas equivalentes.

\subsubsection{¿Qué evalúa?}\label{quuxe9-evaluxfaa-1}

La capacidad para validar diferentes procedimientos algebraicos en la
solución de una ecuación cuadrática.

\subsubsection{Respuesta Correcta: B (Solo María y
Nelson)}\label{respuesta-correcta-b-solo-maruxeda-y-nelson}

\subsubsection{Justificación:}\label{justificaciuxf3n-1}

\begin{itemize}
\tightlist
\item
  María desarrolla correctamente:

  \begin{itemize}
  \tightlist
  \item
    Expande \((x + 2)(x + 3) = 5(x + 3)\) obteniendo
    \(x^2 + 5x + 6 = 5x + 15\)
  \item
    Agrupa términos semejantes: \(x^2 + 6 = 15\)
  \item
    Resuelve para \(x = \pm 3\)
  \end{itemize}
\item
  Nelson aplica factorización:

  \begin{itemize}
  \tightlist
  \item
    Reescribe como \((x + 3)[(x + 2) - 5] = 0\)
  \item
    Obtiene \((x + 3)(x - 3) = 0\)
  \item
    Resuelve para \(x = \pm 3\)
  \end{itemize}
\item
  Óscar comete error al:

  \begin{itemize}
  \tightlist
  \item
    No multiplicar los términos \((x + 2)(x + 3)\)
  \item
    Suma incorrectamente los términos lineales
  \item
    Obtiene solo una solución parcial
  \end{itemize}
\end{itemize}

\subsubsection{Distractores:}\label{distractores-1}

A. Confunde procedimientos válidos con inválidos C. y D. Identifica solo
un procedimiento como correcto cuando hay dos

\subsubsection{Contenidos Matemáticos
Curriculares:}\label{contenidos-matemuxe1ticos-curriculares-2}

\begin{itemize}
\tightlist
\item
  Álgebra y Cálculo
\item
  Período 2

  \begin{itemize}
  \tightlist
  \item
    Álgebra

    \begin{itemize}
    \tightlist
    \item
      Expresiones algebraicas
    \item
      Operaciones con expresiones algebraicas
    \end{itemize}
  \end{itemize}
\end{itemize}

\subsubsection{Tipo: No genérico}\label{tipo-no-genuxe9rico}

\subsubsection{Eje Axial Disciplinar:}\label{eje-axial-disciplinar-1}

Eje 3 - Términos Semejantes, Factorización Básica: Álgebra

\subsubsection{Tarea:}\label{tarea-1}

Evaluar la validez de tres procedimientos diferentes para resolver una
ecuación cuadrática.

\subsubsection{Grado: 9° - 11°}\label{grado-9---11}

\subsection{Competencia 6: Aplicaciones
Integradas}\label{competencia-6-aplicaciones-integradas}

\section{Pregunta 07}\label{pregunta-07}

\begin{center}\rule{0.5\linewidth}{0.5pt}\end{center}

\includegraphics{/home/manjaro_lenovo/Documentos/Math/2023-Matematicas-11-2/07/01.png}

\begin{center}\rule{0.5\linewidth}{0.5pt}\end{center}

\subsection{Análisis}\label{anuxe1lisis-6}

Identifiquemos los elementos clave del problema:

\begin{itemize}
\tightlist
\item
  Total de acciones: 10 + 6 + 4 = 20 acciones
\item
  Ganancias totales: \$80.000.000
\item
  Mario tiene 10 acciones (50\%)
\item
  Carolina tiene 6 acciones (30\%)
\item
  Lucía tiene 4 acciones (20\%)
\end{itemize}

\subsection{Evaluación de
Procedimientos}\label{evaluaciuxf3n-de-procedimientos}

\subsubsection{Procedimiento 1:}\label{procedimiento-1}

Multiplicar ganancias por 10 y dividir entre 100:

\begin{itemize}
\tightlist
\item
  \$80.000.000 × 10 ÷ 100 = \$8.000.000
\end{itemize}

\subsubsection{Procedimiento 2:}\label{procedimiento-2}

Dividir ganancias entre 20 y multiplicar por 10:

\begin{itemize}
\tightlist
\item
  \$80.000.000 ÷ 20 × 10 = \$40.000.000
\end{itemize}

\subsection{Verificación
Matemática}\label{verificaciuxf3n-matemuxe1tica}

Para determinar la parte correcta de Mario (10 acciones de 20 totales):

\begin{itemize}
\tightlist
\item
  Fracción correspondiente: 10/20 = 1/2 = 50\%
\item
  Monto correcto: \$80.000.000 × 50\% = \$40.000.000
\end{itemize}

\subsection{Conclusión}\label{conclusiuxf3n-5}

La respuesta correcta es C: Solo el procedimiento 2 es correcto.

\subsubsection{Justificación:}\label{justificaciuxf3n-2}

\begin{enumerate}
\def\labelenumi{\arabic{enumi}.}
\item
  El procedimiento 1 resulta en \$8.000.000, que es solo el 10\% de las
  ganancias totales, muy por debajo del 50\% que le corresponde a Mario.
\item
  El procedimiento 2 resulta en \$40.000.000, que es exactamente el 50\%
  de las ganancias totales, correspondiente a las 10 acciones de Mario
  de un total de 20.
\end{enumerate}

\subsection{Verificación Final}\label{verificaciuxf3n-final}

\begin{itemize}
\tightlist
\item
  Procedimiento 2: \$80.000.000 ÷ 20 × 10 = \$40.000.000
\item
  Este resultado coincide con el cálculo directo del 50\% de
  \$80.000.000
\item
  Confirma la proporcionalidad requerida según el número de acciones
\end{itemize}

Por lo tanto, solo el procedimiento 2 proporciona el resultado correcto
para calcular la parte correspondiente a Mario.

\begin{center}\rule{0.5\linewidth}{0.5pt}\end{center}

\subsubsection{Análisis de la Pregunta 2023-Matematicas-11-2-07, según
SABER 11
ICFES}\label{anuxe1lisis-de-la-pregunta-2023-matematicas-11-2-07-seguxfan-saber-11-icfes}

\begin{itemize}
\item
  Nivel de Desempeño: 3
\item
  Competencia: Formulación y Ejecución

  \begin{itemize}
  \tightlist
  \item
    Aprendizaje: Frente a un problema que involucre información
    cuantitativa, plantea e implementa estrategias que lleven a
    soluciones adecuadas.
  \item
    Evidencia: Ejecuta un plan de solución para un problema que
    involucra información cuantitativa o esquemática.
  \end{itemize}
\item
  Componente: Numérico-Variacional
\item
  Estándar Asociado: Resuelvo y formulo problemas en situaciones de
  proporcionalidad directa, inversa y producto de medidas.
\item
  ¿Qué evalúa?: La capacidad para validar procedimientos matemáticos que
  involucran proporcionalidad directa en un contexto financiero.
\item
  Respuesta Correcta: C
\item
  Justificación: El procedimiento 2 es correcto porque:

  \(\text{Ganancia de Mario} = \frac{\text{Ganancias totales}}{\text{Total acciones}} \times \text{Acciones de Mario}\)

  \(\text{Ganancia de Mario} = \frac{\$80.000.000}{20} \times 10 = \$40.000.000\)
\item
  Motivos de opciones incorrectas:

  \begin{itemize}
  \tightlist
  \item
    A: Asume que ambos procedimientos dan el mismo resultado
  \item
    B: El procedimiento 1 calcula el 10\% en lugar del 50\% requerido
  \item
    D: No reconoce que el procedimiento 2 es matemáticamente válido
  \end{itemize}
\item
  Contenidos Matemáticos Curriculares:

  \begin{itemize}
  \tightlist
  \item
    Álgebra y Cálculo
  \item
    Período 1

    \begin{itemize}
    \tightlist
    \item
      Números y Operaciones

      \begin{itemize}
      \tightlist
      \item
        Números racionales (fracciones, decimales y porcentajes)
      \end{itemize}
    \end{itemize}
  \end{itemize}
\item
  Genérico: Sí
\item
  Eje Axial Disciplinar: Eje 3 - Función Lineal, Generalizaciones y
  Grado 2
\item
  Tarea: Validar procedimientos aritméticos para calcular una proporción
  directa en un contexto de repartición de ganancias.
\item
  Grado: 7° - 11°
\end{itemize}

\section{Pregunta 08}\label{pregunta-08}

\begin{center}\rule{0.5\linewidth}{0.5pt}\end{center}

\includegraphics{/home/manjaro_lenovo/Documentos/Math/2023-Matematicas-11-2/08/01.png}

\begin{center}\rule{0.5\linewidth}{0.5pt}\end{center}

\subsection{Análisis}\label{anuxe1lisis-7}

Para encontrar la nota necesaria en la cuarta evaluación, debemos: *
Identificar el promedio objetivo: 80 * Reconocer las notas existentes:
70, 93, 89 * Comprender que necesitamos que el promedio de las 4 notas
sea 80

\subsection{Desarrollo Estratégico}\label{desarrollo-estratuxe9gico}

\subsubsection{1. Planteamiento de la
Ecuación}\label{planteamiento-de-la-ecuaciuxf3n}

Para un promedio de 80, sabemos que:

\begin{verbatim}
(70 + 93 + 89 + x) ÷ 4 = 80
\end{verbatim}

donde x es la nota que necesitamos encontrar.

\subsubsection{2. Resolución Algebraica}\label{resoluciuxf3n-algebraica}

\begin{verbatim}
70 + 93 + 89 + x = 80 × 4
70 + 93 + 89 + x = 320
               x = 320 - (70 + 93 + 89)
               x = 320 - 252
               x = 68
\end{verbatim}

\subsubsection{3. Verificación}\label{verificaciuxf3n-1}

Comprobemos que 68 es la respuesta correcta:

\begin{verbatim}
(70 + 93 + 89 + 68) ÷ 4 = 320 ÷ 4 = 80
\end{verbatim}

\subsection{Conclusión}\label{conclusiuxf3n-6}

La respuesta correcta es D. 68

\subsection{Justificación}\label{justificaciuxf3n-3}

\begin{itemize}
\tightlist
\item
  Carlos necesita obtener un 68 en su cuarta evaluación para alcanzar
  exactamente un promedio de 80.
\item
  Este resultado es lógico porque:

  \begin{itemize}
  \tightlist
  \item
    Es una nota posible dentro del rango (1-100)
  \item
    Es menor que sus últimas dos notas (93 y 89)
  \item
    Compensa la nota más baja (70) para lograr el promedio deseado
  \end{itemize}
\end{itemize}

\subsection{Recomendación
Pedagógica}\label{recomendaciuxf3n-pedaguxf3gica}

Este ejercicio demuestra la importancia de:

\begin{enumerate}
\def\labelenumi{\arabic{enumi}.}
\tightlist
\item
  Traducir el problema a lenguaje algebraico
\item
  Resolver paso a paso la ecuación
\item
  Verificar que la solución cumple con todas las condiciones del
  problema
\end{enumerate}

\begin{center}\rule{0.5\linewidth}{0.5pt}\end{center}

\subsection{Análisis de la Pregunta 2023-Matematicas-11-2-08, según
SABER 11
ICFES}\label{anuxe1lisis-de-la-pregunta-2023-matematicas-11-2-08-seguxfan-saber-11-icfes}

\begin{itemize}
\item
  Nivel de Desempeño: 3
\item
  Competencia: Formulación y Ejecución

  \begin{itemize}
  \tightlist
  \item
    Aprendizaje: Frente a un problema que involucre información
    cuantitativa, plantea e implementa estrategias que lleven a
    soluciones adecuadas
  \item
    Evidencia: Ejecuta un plan de solución para un problema que
    involucra información cuantitativa o esquemática
  \end{itemize}
\item
  Componente: Numérico-Variacional
\item
  Estándar Asociado: Resuelvo problemas y simplifico cálculos usando
  propiedades y relaciones de los números reales y de las relaciones y
  operaciones entre ellos
\item
  ¿Qué evalúa?: La capacidad para plantear y resolver una ecuación
  lineal que modele una situación de promedio aritmético
\item
  Respuesta Correcta: D (68)
\item
  Justificación:

  \begin{itemize}
  \tightlist
  \item
    Sea \(x\) la nota de la cuarta evaluación
  \item
    El promedio de 80 se expresa como:
    \(\frac{70 + 93 + 89 + x}{4} = 80\)
  \item
    Resolviendo: \(70 + 93 + 89 + x = 320\)
  \item
    Por lo tanto: \(x = 320 - 252 = 68\)
  \end{itemize}
\item
  Distractores:

  \begin{itemize}
  \tightlist
  \item
    A (63): Resultado de un cálculo incorrecto al no multiplicar 80 por
    4
  \item
    B (43): Error al plantear la ecuación considerando solo tres notas
  \item
    C (84): Confusión al asumir que la nota debe ser mayor que el
    promedio deseado
  \end{itemize}
\item
  Contenidos Matemáticos Curriculares:

  \begin{itemize}
  \tightlist
  \item
    Álgebra y Cálculo
  \item
    Período 1

    \begin{itemize}
    \tightlist
    \item
      Números y Operaciones

      \begin{itemize}
      \tightlist
      \item
        Operaciones aritméticas
      \end{itemize}
    \end{itemize}
  \end{itemize}
\item
  Genérico: Sí
\item
  Eje Axial Disciplinar: Eje 3 - Función Lineal
\item
  Tarea: Determinar un valor desconocido usando el concepto de promedio
  aritmético mediante una ecuación lineal
\item
  Grado: 8° - 11°
\end{itemize}

\section{Pregunta 09}\label{pregunta-09}

\begin{center}\rule{0.5\linewidth}{0.5pt}\end{center}

\includegraphics{/home/manjaro_lenovo/Documentos/Math/2023-Matematicas-11-2/09/01.png}

\begin{center}\rule{0.5\linewidth}{0.5pt}\end{center}

\section{Análisis detallado del
problema}\label{anuxe1lisis-detallado-del-problema}

\subsection{Enunciado}\label{enunciado}

Un auto viaja con rapidez constante. A las \textbf{8:00 a. m.} se
encuentra a \textbf{20 km} del pueblo de partida, y a las \textbf{10:00
a. m.} está a \textbf{100 km} del mismo pueblo. Se busca calcular la
rapidez media del auto en km/h a partir de tres pasos presentados, y
determinar si hay errores en el procedimiento.

\subsubsection{Pasos presentados en el
problema}\label{pasos-presentados-en-el-problema}

\begin{enumerate}
\def\labelenumi{\arabic{enumi}.}
\item
  \textbf{Paso 1:} Calcular la distancia recorrida:\\
  \(100 \, \text{km} - 20 \, \text{km} = 80 \, \text{km}\).
\item
  \textbf{Paso 2:} Calcular el tiempo transcurrido:\\
  \(10 \, \text{h} - 8 \, \text{h} = 2 \, \text{h}\).
\item
  \textbf{Paso 3:} Dividir los valores obtenidos para hallar la rapidez
  media:\\
  \(\frac{80 \, \text{km}}{2 \, \text{h}} = 40 \, \text{km/h}\).
\end{enumerate}

\subsection{Alternativas de respuesta}\label{alternativas-de-respuesta}

\begin{itemize}
\tightlist
\item
  \textbf{A.} En el paso 1, ya que la distancia recorrida es
  \(100 \, \text{km}\) y no \(80 \, \text{km}\).\\
\item
  \textbf{B.} En el paso 2, ya que el tiempo transcurrido es
  \(8 \, \text{h} - 10 \, \text{h} = -2 \, \text{h}\).\\
\item
  \textbf{C.} En el paso 3, ya que el cociente es
  \(80 \, \text{km}/2 \, \text{h} = 40 \, \text{h/km}\).\\
\item
  \textbf{D.} En el paso 1, ya que la distancia recorrida es
  \(20 \, \text{km} - 100 \, \text{km} = -20 \, \text{km}\).
\end{itemize}

\subsection{Solución paso a paso}\label{soluciuxf3n-paso-a-paso}

\begin{enumerate}
\def\labelenumi{\arabic{enumi}.}
\item
  \textbf{Paso 1: Distancia recorrida}\\
  La distancia recorrida por el auto se calcula correctamente como:\\
  \[
  \text{Distancia} = 100 \, \text{km} - 20 \, \text{km} = 80 \, \text{km}.
  \]\\
  No hay error en este paso. El resultado es correcto y coincide con los
  datos del problema.
\item
  \textbf{Paso 2: Tiempo transcurrido}\\
  El tiempo transcurrido también se calcula correctamente:\\
  \[
  \text{Tiempo} = 10 \, \text{h} - 8 \, \text{h} = 2 \, \text{h}.
  \]\\
  El valor obtenido es correcto y corresponde al intervalo entre las
  \textbf{8:00 a. m.} y las \textbf{10:00 a. m.}.
\item
  \textbf{Paso 3: Cálculo de la rapidez media}\\
  La rapidez media se calcula dividiendo la distancia recorrida entre el
  tiempo transcurrido:\\
  \[
  \text{Rapidez media} = \frac{80 \, \text{km}}{2 \, \text{h}} = 40 \, \text{km/h}.
  \]\\
  El cálculo es correcto, pero hay un error conceptual en la unidad de
  medida que se describe en el problema: se menciona \(\text{h/km}\) en
  lugar de \(\text{km/h}\). Sin embargo, este error es menor y no afecta
  el resultado matemático.
\end{enumerate}

\begin{center}\rule{0.5\linewidth}{0.5pt}\end{center}

\subsection{Conclusión}\label{conclusiuxf3n-7}

La \textbf{opción correcta} es: \textbf{C}.\\
El error en el paso 3 radica en la descripción de la unidad como
\(\text{h/km}\), lo cual es incorrecto; debería ser \(\text{km/h}\).
Esto refleja una imprecisión en la presentación del resultado final,
aunque los cálculos matemáticos sean correctos.

\begin{center}\rule{0.5\linewidth}{0.5pt}\end{center}

\subsection{Notas adicionales}\label{notas-adicionales}

La respuesta final, si se redactara correctamente, sería:\\
\[
\text{Rapidez media} = 40 \, \text{km/h}.
\]

\begin{center}\rule{0.5\linewidth}{0.5pt}\end{center}

\subsection{Análisis de la Pregunta 2023-Matematicas-11-2-09, según
SABER 11
ICFES}\label{anuxe1lisis-de-la-pregunta-2023-matematicas-11-2-09-seguxfan-saber-11-icfes}

\subsection{Nivel de Desempeño}\label{nivel-de-desempeuxf1o-1}

3

\subsection{Competencia}\label{competencia-1}

\begin{itemize}
\tightlist
\item
  Argumentación

  \begin{itemize}
  \tightlist
  \item
    Afirmación: Valida procedimientos y estrategias matemáticas
    utilizadas para dar solución a problemas
  \item
    Evidencia: Argumenta a favor o en contra de un procedimiento para
    resolver un problema a la luz de criterios presentados o
    establecidos
  \end{itemize}
\end{itemize}

\subsection{Componente}\label{componente-1}

Numérico-Variacional

\subsection{Estándar Asociado}\label{estuxe1ndar-asociado-2}

Utilizo las técnicas de aproximación en procesos infinitos numéricos.

\subsection{¿Qué evalúa?}\label{quuxe9-evaluxfaa-2}

La capacidad del estudiante para identificar errores en procedimientos
que involucran el cálculo de la rapidez media, incluyendo el manejo
correcto de unidades físicas.

\subsection{Respuesta Correcta}\label{respuesta-correcta}

C. En el paso 3, ya que el cociente es \(80\) km/\(2\) h = \(40\) h/km

\subsection{Justificación}\label{justificaciuxf3n-4}

El error está en el paso 3 porque:

\begin{enumerate}
\def\labelenumi{\arabic{enumi}.}
\tightlist
\item
  La unidad de medida de la rapidez media debe ser km/h (kilómetros por
  hora)
\item
  En el procedimiento se expresa incorrectamente como h/km (horas por
  kilómetro)
\item
  Aunque el valor numérico 40 es correcto, la unidad está invertida
\end{enumerate}

\subsection{Distractores}\label{distractores-2}

\begin{itemize}
\tightlist
\item
  A: Confunde el desplazamiento total con la posición final
\item
  B: Invierte incorrectamente el orden de la resta en el tiempo
\item
  D: Invierte incorrectamente el orden de la resta en la distancia
\end{itemize}

\subsection{Contenidos Matemáticos
Curriculares}\label{contenidos-matemuxe1ticos-curriculares-3}

\begin{itemize}
\tightlist
\item
  Álgebra y Cálculo
\item
  Período 1

  \begin{itemize}
  \tightlist
  \item
    Números y Operaciones

    \begin{itemize}
    \tightlist
    \item
      Uso de números y operaciones en la resolución de problemas
    \end{itemize}
  \end{itemize}
\end{itemize}

\subsection{¿Genérico o No
genérico?}\label{genuxe9rico-o-no-genuxe9rico-1}

Genérico

\subsection{Eje Axial Disciplinar}\label{eje-axial-disciplinar-2}

Eje 3: Funciones I - Función Lineal

\subsection{Tarea}\label{tarea-2}

Validar la coherencia entre las unidades de medida y el significado
físico de la rapidez media en un problema de movimiento rectilíneo
uniforme.

\subsection{Grado}\label{grado-1}

11°

\section{Pregunta 10}\label{pregunta-10}

\begin{center}\rule{0.5\linewidth}{0.5pt}\end{center}

\includegraphics{/home/manjaro_lenovo/Documentos/Math/2023-Matematicas-11-2/10/01.png}

\begin{center}\rule{0.5\linewidth}{0.5pt}\end{center}

\section{Análisis detallado del
problema}\label{anuxe1lisis-detallado-del-problema-1}

\subsubsection{Procedimiento Propuesto}\label{procedimiento-propuesto}

\begin{enumerate}
\def\labelenumi{\arabic{enumi}.}
\tightlist
\item
  \textbf{Paso 1:} Sumar las áreas de un bloque grande, uno mediano y
  uno pequeño, lo que equivale a 9.000 cm².
\item
  \textbf{Paso 2:} Dividir 20.000 cm² entre 9.000 cm².
\end{enumerate}

\subsubsection{Análisis del
Procedimiento}\label{anuxe1lisis-del-procedimiento}

\begin{enumerate}
\def\labelenumi{\arabic{enumi}.}
\item
  \textbf{Paso 1:}

  \begin{itemize}
  \tightlist
  \item
    Área de un bloque grande: 6.500 cm²
  \item
    Área de un bloque mediano: 2.000 cm²
  \item
    Área de un bloque pequeño: 500 cm²
  \end{itemize}

  Sumando estas áreas: \[
  6.500 \, \text{cm}^2 + 2.000 \, \text{cm}^2 + 500 \, \text{cm}^2 = 9.000 \, \text{cm}^2
  \]
\item
  \textbf{Paso 2:}

  \begin{itemize}
  \tightlist
  \item
    Dividir el área total del muro (20.000 cm²) entre el área total de
    los tres bloques (9.000 cm²): \[
    \frac{20.000 \, \text{cm}^2}{9.000 \, \text{cm}^2} \approx 2.22
    \]
  \end{itemize}

  Este resultado no es un número entero, lo que indica que no se puede
  cubrir exactamente el muro con bloques enteros de estos tamaños sin
  dejar espacios o sin cortar bloques.
\end{enumerate}

\subsubsection{Evaluación de las
Opciones}\label{evaluaciuxf3n-de-las-opciones}

\begin{itemize}
\tightlist
\item
  \textbf{A.} Sí, pues solo son necesarias las medidas totales de ambas
  áreas, la del muro y la de los bloques.

  \begin{itemize}
  \tightlist
  \item
    \textbf{Evaluación:} Esta opción es incorrecta porque, aunque se
    usan las medidas totales, el resultado final no es un número entero,
    lo que hace que el procedimiento no sea válido.
  \end{itemize}
\item
  \textbf{B.} No, pues al hacer la división no se consideran las
  dimensiones del muro y de los bloques.

  \begin{itemize}
  \tightlist
  \item
    \textbf{Evaluación:} Esta opción es correcta. El procedimiento no
    considera las dimensiones específicas del muro y de los bloques, lo
    que es crucial para determinar si los bloques pueden cubrir el muro
    sin dejar espacios.
  \end{itemize}
\item
  \textbf{C.} Sí, pues se asegura el uso de cada material, por lo menos,
  una vez.

  \begin{itemize}
  \tightlist
  \item
    \textbf{Evaluación:} Esta opción es incorrecta porque, aunque se use
    cada material al menos una vez, el resultado final no es un número
    entero, lo que hace que el procedimiento no sea válido.
  \end{itemize}
\item
  \textbf{D.} No, pues el resultado final, al realizar la división, no
  es un número entero.

  \begin{itemize}
  \tightlist
  \item
    \textbf{Evaluación:} Esta opción es correcta. El resultado final de
    la división no es un número entero, lo que indica que el
    procedimiento no es válido para determinar la cantidad mínima de
    bloques necesarios.
  \end{itemize}
\end{itemize}

\subsubsection{Conclusión}\label{conclusiuxf3n-8}

Las opciones correctas son \textbf{B} y \textbf{D}. Sin embargo, la
opción \textbf{B} es más completa porque explica que el procedimiento no
considera las dimensiones específicas del muro y de los bloques, lo que
es crucial para determinar si los bloques pueden cubrir el muro sin
dejar espacios.

Por lo tanto, la opción correcta es:

\textbf{B. No, pues al hacer la división no se consideran las
dimensiones del muro y de los bloques.}

\begin{center}\rule{0.5\linewidth}{0.5pt}\end{center}

\subsection{Análisis de la Pregunta 2023-Matematicas-11-2-10, según
SABER 11
ICFES}\label{anuxe1lisis-de-la-pregunta-2023-matematicas-11-2-10-seguxfan-saber-11-icfes}

\subsection{Nivel, Competencia y
Componentes}\label{nivel-competencia-y-componentes}

\begin{itemize}
\tightlist
\item
  \textbf{Nivel de Desempeño}: 4
\item
  \textbf{Competencia}: Argumentación

  \begin{itemize}
  \tightlist
  \item
    \textbf{Afirmación}: Valida procedimientos y estrategias matemáticas
    utilizadas para dar solución a problemas
  \item
    \textbf{Evidencia}: Argumenta a favor o en contra de un
    procedimiento para resolver un problema a la luz de criterios
    presentados o establecidos
  \end{itemize}
\item
  \textbf{Componente}: Numérico-Variacional
\item
  \textbf{Estándar Asociado}: Justifico la pertinencia de utilizar
  unidades de medida estandarizadas en situaciones tomadas de distintas
  ciencias.
\end{itemize}

\subsection{Análisis de la Pregunta}\label{anuxe1lisis-de-la-pregunta}

\subsubsection{¿Qué evalúa?}\label{quuxe9-evaluxfaa-3}

Evalúa la capacidad del estudiante para analizar críticamente un
procedimiento propuesto para resolver un problema geométrico que
involucra áreas y dimensiones.

\subsubsection{Respuesta Correcta: B}\label{respuesta-correcta-b}

No, pues al hacer la división no se consideran las dimensiones del muro
y de los bloques.

\subsubsection{Justificación}\label{justificaciuxf3n-5}

El procedimiento es incorrecto porque:

\begin{enumerate}
\def\labelenumi{\arabic{enumi}.}
\tightlist
\item
  Solo considera las áreas totales sin tener en cuenta:

  \begin{itemize}
  \tightlist
  \item
    Las dimensiones específicas del muro (largo y ancho)
  \item
    Las dimensiones de cada tipo de bloque
  \end{itemize}
\item
  La división \(\frac{20.000}{9.000} \approx 2.22\) no garantiza que los
  bloques puedan cubrir el muro de manera exacta sin dejar espacios.
\end{enumerate}

\subsubsection{Contenidos Matemáticos
Curriculares}\label{contenidos-matemuxe1ticos-curriculares-4}

\begin{itemize}
\tightlist
\item
  \textbf{Categoría}: Geometría
\item
  \textbf{Período}: 2

  \begin{itemize}
  \tightlist
  \item
    Polígonos

    \begin{itemize}
    \tightlist
    \item
      Área y perímetro
    \item
      Áreas por composición y descomposición
    \end{itemize}
  \end{itemize}
\end{itemize}

\subsubsection{Características
Adicionales}\label{caracteruxedsticas-adicionales}

\begin{itemize}
\tightlist
\item
  \textbf{Tipo}: No genérico
\item
  \textbf{Eje Axial Disciplinar}: Área, Perímetro, Subconjuntos de Áreas
\item
  \textbf{Tarea}: Evaluar la validez de un procedimiento para determinar
  la cantidad mínima de bloques necesarios para cubrir un área
  rectangular dada
\item
  \textbf{Grado sugerido}: 9° - 11°
\end{itemize}

\section{Pregunta 11}\label{pregunta-11}

\begin{center}\rule{0.5\linewidth}{0.5pt}\end{center}

\includegraphics{/home/manjaro_lenovo/Documentos/Math/2023-Matematicas-11-2/11/01.png}

\begin{center}\rule{0.5\linewidth}{0.5pt}\end{center}

\section{Análisis detallado del
problema}\label{anuxe1lisis-detallado-del-problema-2}

La respuesta correcta es la \textbf{B}.

\textbf{Explicación detallada:}

El estudiante se equivoca al pensar que todo número racional entre
\(\frac{1}{7}\) y \(\frac{1}{5}\) debe tener denominador 6. Para
comprobar esto, podemos buscar un contraejemplo, es decir, un número
racional entre \(\frac{1}{7}\) y \(\frac{1}{5}\) que no tenga
denominador 6.

\textbf{Paso 1: Comparar las fracciones dadas.}

Para comparar las fracciones y encontrar números entre ellas, es útil
tener un denominador común. El mínimo común múltiplo de 7 y 5 es 35.

\(\frac{1}{7} = \frac{1 \times 5}{7 \times 5} = \frac{5}{35}\)

\(\frac{1}{5} = \frac{1 \times 7}{5 \times 7} = \frac{7}{35}\)

Ahora vemos que los números racionales que buscamos están entre
\(\frac{5}{35}\) y \(\frac{7}{35}\).

\textbf{Paso 2: Analizar las opciones.}

\begin{itemize}
\item
  \textbf{Opción A:} Afirma que \(\frac{1}{6}\) está entre los dos
  números. Vamos a comprobarlo. Para comparar \(\frac{1}{6}\) con
  \(\frac{1}{7}\) y \(\frac{1}{5}\), también podemos usar un denominador
  común. El mínimo común múltiplo de 7, 5 y 6 es 210.
  \(\frac{1}{7} = \frac{30}{210}\) \(\frac{1}{6} = \frac{35}{210}\)
  \(\frac{1}{5} = \frac{42}{210}\) Como \(30 < 35 < 42\), entonces
  \(\frac{1}{7} < \frac{1}{6} < \frac{1}{5}\). Si bien \(\frac{1}{6}\)
  está entre los dos números, esto no demuestra que \emph{todos} los
  números entre ellos deban tener denominador 6.
\item
  \textbf{Opción B:} Afirma que \(\frac{12}{70}\) está entre los dos
  números. Vamos a comprobarlo. Primero, simplificamos la fracción
  \(\frac{12}{70}\):
  \(\frac{12}{70} = \frac{12 \div 2}{70 \div 2} = \frac{6}{35}\) Ya
  habíamos expresado \(\frac{1}{7}\) y \(\frac{1}{5}\) con denominador
  35: \(\frac{1}{7} = \frac{5}{35}\) \(\frac{1}{5} = \frac{7}{35}\) Como
  \(5 < 6 < 7\), entonces
  \(\frac{5}{35} < \frac{6}{35} < \frac{7}{35}\), lo que significa que
  \(\frac{1}{7} < \frac{12}{70} < \frac{1}{5}\). La fracción
  \(\frac{12}{70}\) (o su forma simplificada \(\frac{6}{35}\)) está
  entre los dos números y su denominador (70 o 35) no es 6. Esto
  demuestra que la suposición del estudiante es incorrecta.
\item
  \textbf{Opción C:} Indica que cuando el denominador es más grande, la
  fracción es más pequeña y 6 es más grande que 5. Si bien la primera
  parte de la afirmación es correcta (para numeradores positivos
  iguales), la conclusión no justifica la afirmación del estudiante. La
  comparación de los denominadores 5 y 6 no prueba que todos los números
  entre \(\frac{1}{7}\) y \(\frac{1}{5}\) deban tener denominador 6.
\item
  \textbf{Opción D:} Afirma que entre los dos números hay infinitos
  números, así que alguno debe tener denominador 6. Si bien es cierto
  que hay infinitos números racionales entre dos racionales
  cualesquiera, esto no implica que \emph{todos} los números entre ellos
  deban tener un denominador específico como 6. De hecho, la existencia
  de infinitos números sugiere que hay muchas más opciones de
  denominadores.
\end{itemize}

\textbf{Conclusión:}

La opción correcta es la \textbf{B} porque proporciona un contraejemplo
específico: la fracción \(\frac{12}{70}\) está entre \(\frac{1}{7}\) y
\(\frac{1}{5}\), y su denominador no es 6. Esto demuestra que la
suposición del estudiante es incorrecta.

\begin{center}\rule{0.5\linewidth}{0.5pt}\end{center}

\subsection{Analisis de la Pregunta 2023-Matematicas-11-2-11 segun SABER
11
ICFES}\label{analisis-de-la-pregunta-2023-matematicas-11-2-11-segun-saber-11-icfes}

\begin{itemize}
\item
  Nivel de Desempeño: 4
\item
  Competencia: Argumentación

  \begin{itemize}
  \tightlist
  \item
    Afirmación: Valida procedimientos y estrategias matemáticas
    utilizadas para dar solución a problemas
  \item
    Evidencia: Plantea afirmaciones que sustentan o refutan una
    interpretación dada a la información disponible en el marco de la
    solución de un problema
  \end{itemize}
\item
  Componente: Numérico-Variacional
\item
  Estándar Asociado: Justifico la pertinencia de utilizar unidades de
  medida estandarizadas en situaciones tomadas de distintas ciencias
\item
  ¿Qué evalúa?: La capacidad de refutar una afirmación sobre números
  racionales usando un contraejemplo específico
\item
  Respuesta Correcta: B, porque proporciona un contraejemplo
  (\(\frac{12}{70}\)) que demuestra que no todo número racional entre
  \(\frac{1}{7}\) y \(\frac{1}{5}\) debe tener denominador 6
\item
  Distractores:

  \begin{itemize}
  \tightlist
  \item
    A: Confunde la existencia de un número con denominador 6 entre los
    dados con la necesidad de que todos tengan ese denominador
  \item
    C: Usa una propiedad válida de fracciones pero llega a una
    conclusión incorrecta
  \item
    D: Confunde la existencia de infinitos números con la necesidad de
    tener un denominador específico
  \end{itemize}
\item
  Contenidos Matemáticos Curriculares:

  \begin{itemize}
  \tightlist
  \item
    Álgebra y Cálculo
  \item
    Período 1

    \begin{itemize}
    \tightlist
    \item
      Números y Operaciones

      \begin{itemize}
      \tightlist
      \item
        Números racionales (fracciones, decimales y porcentajes)
      \end{itemize}
    \end{itemize}
  \end{itemize}
\item
  No genérico
\item
  Eje Axial Disciplinar: Eje 3 - Teoría de Números
\item
  Tarea: Refutar una afirmación sobre números racionales mediante la
  construcción de un contraejemplo
\item
  Grado: 11° (También aplicable desde 8° según los Estándares Básicos de
  Competencias)
\end{itemize}

\section{Pregunta 12}\label{pregunta-12}

\begin{center}\rule{0.5\linewidth}{0.5pt}\end{center}

\includegraphics{/home/manjaro_lenovo/Documentos/Math/2023-Matematicas-11-2/12/all.png}

\begin{center}\rule{0.5\linewidth}{0.5pt}\end{center}

\section{Análisis detallado del
problema}\label{anuxe1lisis-detallado-del-problema-3}

Vamos a analizar la pregunta y la solución paso a paso.

La desigualdad planteada es: 20.000 + 1.000x \textless{} 15.000 + 1.200x

Esta desigualdad busca encontrar la condición en la que el costo total
de la \textbf{Opción 1} es \textbf{menor que} el costo total de la
\textbf{Opción 2}.

Resolvamos la desigualdad:

\begin{enumerate}
\def\labelenumi{\arabic{enumi}.}
\item
  Restamos 1.000x a ambos lados: 20.000 \textless{} 15.000 + 0.200x
\item
  Restamos 15.000 a ambos lados: 5.000 \textless{} 0.200x
\item
  Dividimos ambos lados por 0.200: 5.000 / 0.200 \textless{} x 25
  \textless{} x
\end{enumerate}

La solución a la desigualdad es x \textgreater{} 25.

¿Qué significa esta solución en el contexto del problema?

\begin{itemize}
\tightlist
\item
  \textbf{x} representa la cantidad de atracciones.
\item
  La desigualdad se planteó para cuando el costo de la Opción 1 es menor
  que el costo de la Opción 2.
\item
  La solución x \textgreater{} 25 significa que para cualquier cantidad
  de atracciones mayor a 25, el costo de la Opción 1 es menor que el
  costo de la Opción 2.
\end{itemize}

Ahora, revisemos las opciones de respuesta:

\begin{itemize}
\item
  \textbf{A. El rango de costos que una persona debería pagar si escoge
  la opción 1.} La desigualdad nos ayuda a determinar cuándo la Opción 1
  es más barata, pero la solución se refiere directamente a la cantidad
  de atracciones, no al costo en sí.
\item
  \textbf{B. El rango de la cantidad de atracciones en las que el costo
  de la opción 1 es mayor que el de la opción 2.} La desigualdad está
  planteada para cuando el costo de la Opción 1 es \emph{menor que} el
  de la Opción 2. La solución refleja esta condición.
\item
  \textbf{C. El rango de la cantidad de atracciones en las que el costo
  de la opción 1 es menor que el de la opción 2.} Esta opción coincide
  con nuestra interpretación de la solución x \textgreater{} 25. Indica
  el rango de cantidad de atracciones para el cual elegir la Opción 1 es
  más económico.
\item
  \textbf{D. El rango de costos que una persona debería pagar si escoge
  la opción 2.} Al igual que la opción A, la desigualdad ayuda a
  determinar cuándo la Opción 1 es más barata, y la solución se refiere
  a la cantidad de atracciones, no al costo de la Opción 2.
\end{itemize}

Por lo tanto, la respuesta correcta es la \textbf{C}. La solución de la
desigualdad nos da el rango de la cantidad de atracciones en las que el
costo de la opción 1 es menor que el de la opción 2.

\begin{center}\rule{0.5\linewidth}{0.5pt}\end{center}

\subsection{Análisis de la Pregunta 2023-Matematicas-11-2-12, según
SABER 11
ICFES}\label{anuxe1lisis-de-la-pregunta-2023-matematicas-11-2-12-seguxfan-saber-11-icfes}

\textbf{Nivel de Desempeño}: 3

\textbf{Competencia}: Interpretación y Representación

\begin{itemize}
\tightlist
\item
  \emph{Afirmación}: Comprende y transforma la información cuantitativa
  y esquemática presentada en distintos formatos
\item
  \emph{Evidencia}: Da cuenta de las características básicas de la
  información presentada en diferentes formatos
\end{itemize}

\textbf{Componente}: Numérico-Variacional

\textbf{Estándar Asociado}: Resuelvo y formulo problemas que involucren
relaciones y propiedades de semejanza y congruencia usando
representaciones visuales.

\textbf{¿Qué evalúa?}: La capacidad de interpretar y resolver una
desigualdad lineal en un contexto de costos, identificando el
significado de la solución.

\textbf{Respuesta Correcta}: C. El rango de la cantidad de atracciones
en las que el costo de la opción 1 es menor que el de la opción 2.

\textbf{Justificación}: La desigualdad
\(20.000 + 1.000x < 15.000 + 1.200x\) se resuelve:

\begin{enumerate}
\def\labelenumi{\arabic{enumi}.}
\tightlist
\item
  \(20.000 - 15.000 < 1.200x - 1.000x\)
\item
  \(5.000 < 200x\)
\item
  \(x > 25\)
\end{enumerate}

Esto significa que para más de 25 atracciones, la opción 1 resulta más
económica.

\textbf{Distractores}:

\begin{itemize}
\tightlist
\item
  A: Confunde el resultado con costos en lugar de cantidad de
  atracciones
\item
  B: Interpreta incorrectamente el símbolo de desigualdad
\item
  D: Similar a A, confunde variables de la solución
\end{itemize}

\textbf{Contenidos Matemáticos Curriculares}:

\begin{itemize}
\tightlist
\item
  \emph{Álgebra y Cálculo}
\item
  \emph{Período}: 2

  \begin{itemize}
  \tightlist
  \item
    Expresiones algebraicas
  \item
    Operaciones con expresiones algebraicas
  \end{itemize}
\end{itemize}

\textbf{Tipo}: No genérico

\textbf{Eje Axial Disciplinar}: Función Lineal, Generalizaciones y Grado
2: Funciones I

\textbf{Tarea}: Resolver una desigualdad lineal e interpretar su
solución en el contexto de un problema de costos.

\textbf{Grado sugerido}: 9° - 11°

\section{Pregunta 13}\label{pregunta-13}

\begin{center}\rule{0.5\linewidth}{0.5pt}\end{center}

\includegraphics{/home/manjaro_lenovo/Documentos/Math/2023-Matematicas-11-2/13/all.png}

\begin{center}\rule{0.5\linewidth}{0.5pt}\end{center}

\subsection{Análisis detallado del
problema}\label{anuxe1lisis-detallado-del-problema-4}

\subsubsection{\texorpdfstring{Paso 1: Sumar los porcentajes de
participación de las empresas \(T_1\) y
\(T_2\)}{Paso 1: Sumar los porcentajes de participación de las empresas T\_1 y T\_2}}\label{paso-1-sumar-los-porcentajes-de-participaciuxf3n-de-las-empresas-t_1-y-t_2}

De la gráfica, tenemos: - \(T_1 = 58\%\) - \(T_2 = 22\%\)

Sumamos estos porcentajes: \[ T_1 + T_2 = 58\% + 22\% = 80\% \]

\subsubsection{Paso 2: Multiplicar el resultado del paso 1 por el total
de usuarios de telefonía móvil en el
país}\label{paso-2-multiplicar-el-resultado-del-paso-1-por-el-total-de-usuarios-de-telefonuxeda-muxf3vil-en-el-pauxeds}

Sea \(U\) el total de usuarios de telefonía móvil en el país.
Multiplicamos el resultado del paso 1 por \(U\): \[ 0.80 \times U \]

\subsubsection{Paso 3: Dividir el resultado del paso 2 entre
100}\label{paso-3-dividir-el-resultado-del-paso-2-entre-100}

Dividimos el resultado del paso 2 entre 100:
\[ \frac{0.80 \times U}{100} = 0.008 \times U \]

Este resultado, \(0.008 \times U\), representa el número de usuarios que
usan una empresa de telefonía móvil extranjera.

\subsubsection{Análisis de las
opciones}\label{anuxe1lisis-de-las-opciones}

\begin{itemize}
\tightlist
\item
  \textbf{A. El número de usuarios que usa una empresa de telefonía
  móvil nacional.}

  \begin{itemize}
  \tightlist
  \item
    No se puede determinar con el procedimiento dado, ya que no se suman
    los porcentajes de \(T_3\) y \(T_4\).
  \end{itemize}
\item
  \textbf{B. El número de usuarios que usa una empresa de telefonía
  móvil extranjera.}

  \begin{itemize}
  \tightlist
  \item
    Sí se puede determinar con el procedimiento dado, como se mostró en
    los pasos anteriores.
  \end{itemize}
\item
  \textbf{C. El porcentaje total de usuarios que usa una empresa de
  telefonía móvil nacional.}

  \begin{itemize}
  \tightlist
  \item
    No se puede determinar con el procedimiento dado, ya que no se suman
    los porcentajes de \(T_3\) y \(T_4\).
  \end{itemize}
\item
  \textbf{D. El porcentaje total de usuarios que usa una empresa de
  telefonía móvil extranjera.}

  \begin{itemize}
  \tightlist
  \item
    No se puede determinar directamente con el procedimiento dado, ya
    que el resultado obtenido es el número de usuarios, no el
    porcentaje.
  \end{itemize}
\end{itemize}

\subsubsection{Conclusión}\label{conclusiuxf3n-9}

La opción correcta es:

\textbf{B. El número de usuarios que usa una empresa de telefonía móvil
extranjera.}

Este es el dato que se puede determinar mediante el procedimiento
propuesto.

\begin{center}\rule{0.5\linewidth}{0.5pt}\end{center}

\subsection{Análisis de la Pregunta 2023-Matematicas-11-2-13, según
SABER 11
ICFES}\label{anuxe1lisis-de-la-pregunta-2023-matematicas-11-2-13-seguxfan-saber-11-icfes}

\subsection{Nivel de Desempeño: 3}\label{nivel-de-desempeuxf1o-3-1}

\subsection{Competencia: Interpretación y
Representación}\label{competencia-interpretaciuxf3n-y-representaciuxf3n}

\begin{itemize}
\tightlist
\item
  \textbf{Afirmación}: Comprende y transforma la información
  cuantitativa y esquemática presentada en distintos formatos
\item
  \textbf{Evidencia}: Da cuenta de las características básicas de la
  información presentada en diferentes formatos como series, gráficas,
  tablas y esquemas
\end{itemize}

\subsection{Componente:
Numérico-Variacional}\label{componente-numuxe9rico-variacional-1}

\subsection{Estándar Asociado:}\label{estuxe1ndar-asociado-3}

Interpreto y comparo resultados de estudios con información estadística
provenientes de medios de comunicación.

\subsection{¿Qué evalúa?}\label{quuxe9-evaluxfaa-4}

La capacidad del estudiante para interpretar información presentada en
un gráfico de barras y realizar operaciones aritméticas básicas con
porcentajes para determinar valores específicos.

\subsection{Respuesta Correcta: B}\label{respuesta-correcta-b-1}

\subsubsection{Justificación:}\label{justificaciuxf3n-6}

El procedimiento descrito permite calcular el número de usuarios de
telefonía móvil extranjera:

\begin{enumerate}
\def\labelenumi{\arabic{enumi}.}
\tightlist
\item
  \(T_1 + T_2 = 58\% + 22\% = 80\%\)
\item
  Si \(U\) es el total de usuarios: \(80\% \times U\)
\item
  \(\frac{80\% \times U}{100} = 0.8 \times U\)
\end{enumerate}

Este resultado representa el número total de usuarios que utilizan
empresas extranjeras.

\subsection{Motivos de las Opciones
Incorrectas:}\label{motivos-de-las-opciones-incorrectas-1}

\begin{itemize}
\tightlist
\item
  \textbf{A}: El procedimiento no suma los porcentajes de empresas
  nacionales (\(T_3\) y \(T_4\))
\item
  \textbf{C}: El procedimiento no calcula porcentajes de empresas
  nacionales
\item
  \textbf{D}: El resultado es una cantidad absoluta, no un porcentaje
\end{itemize}

\subsection{Contenidos Matemáticos
Curriculares:}\label{contenidos-matemuxe1ticos-curriculares-5}

\begin{itemize}
\tightlist
\item
  \textbf{Álgebra y Cálculo}
\item
  \textbf{Período}: 1

  \begin{itemize}
  \tightlist
  \item
    Números y Operaciones

    \begin{itemize}
    \tightlist
    \item
      Números racionales (porcentajes)
    \item
      Operaciones aritméticas
    \end{itemize}
  \end{itemize}
\end{itemize}

\subsection{Tipo: Genérico}\label{tipo-genuxe9rico}

\subsection{Eje Axial Disciplinar:}\label{eje-axial-disciplinar-3}

Eje 1 - Univariadas: Interpretación de tablas y barras elementales

\subsection{Tarea:}\label{tarea-3}

Interpretar un gráfico de barras y realizar operaciones con porcentajes
para determinar el número de usuarios de telefonía móvil extranjera.

\subsection{Grado: 11°}\label{grado-11}

\section{Pregunta 14}\label{pregunta-14}

\begin{center}\rule{0.5\linewidth}{0.5pt}\end{center}

\includegraphics{/home/manjaro_lenovo/Documentos/Math/2023-Matematicas-11-2/14/all.png}

\begin{center}\rule{0.5\linewidth}{0.5pt}\end{center}

\subsection{Análisis detallado del
problema}\label{anuxe1lisis-detallado-del-problema-5}

Vamos a simplificar la expresión \(-(8a + 2b - 5a - 4b)4a\) paso a paso
y verificar cuál de las opciones (A, B, C o D) es la correcta.

\subsection{Paso 1: Agrupar términos
semejantes}\label{paso-1-agrupar-tuxe9rminos-semejantes}

Primero, agrupamos los términos semejantes dentro del paréntesis:

\[
-(8a - 5a + 2b - 4b)4a
\]

Simplificamos los términos dentro del paréntesis:

\[
-(3a - 2b)4a
\]

\subsection{Paso 2: Operar y cambiar los signos de los términos en
paréntesis}\label{paso-2-operar-y-cambiar-los-signos-de-los-tuxe9rminos-en-paruxe9ntesis}

Ahora, operamos el signo menos fuera del paréntesis, lo que cambia los
signos de los términos dentro del paréntesis:

\[
(3a - 2b)4a
\]

\subsection{Paso 3: Multiplicar los
factores}\label{paso-3-multiplicar-los-factores}

Multiplicamos los factores:

\[
3a \cdot 4a - 2b \cdot 4a
\]

Esto se simplifica a:

\[
12a^2 - 8ab
\]

\subsection{Paso 4: Factorizar}\label{paso-4-factorizar}

Finalmente, factorizamos la expresión resultante:

\[
4(3a^2 - 2ab)
\]

\subsection{Análisis de las
opciones}\label{anuxe1lisis-de-las-opciones-1}

\begin{itemize}
\tightlist
\item
  \textbf{Opción A}: En el paso 1, pues debió sumar todos los términos,
  así: \(-(8a - 5a + 2b - 4b)4a = -(a)4a\).

  \begin{itemize}
  \tightlist
  \item
    Esta opción es incorrecta porque no se sumaron correctamente los
    términos dentro del paréntesis.
  \end{itemize}
\item
  \textbf{Opción B}: En el 2, pues no consideró que el signo menos
  cambia los signos dentro del paréntesis, así: \((-3a - 2b)4a\).

  \begin{itemize}
  \tightlist
  \item
    Esta opción es incorrecta porque no se cambiaron correctamente los
    signos dentro del paréntesis.
  \end{itemize}
\item
  \textbf{Opción C}: En el 3, pues no tuvo en cuenta que el segundo
  término tiene una variable diferente, así: \(-12a^2 + 8ab\).

  \begin{itemize}
  \tightlist
  \item
    Esta opción es incorrecta porque no se tuvo en cuenta la variable
    diferente en el segundo término.
  \end{itemize}
\item
  \textbf{Opción D}: En el 4, pues el factor común entre ambos términos
  es 4 y debió quedar así: \(4(-3a^2 + 2a)\).

  \begin{itemize}
  \tightlist
  \item
    Esta opción es incorrecta porque el factor común y la factorización
    no se realizaron correctamente.
  \end{itemize}
\end{itemize}

\subsection{Conclusión}\label{conclusiuxf3n-10}

La opción correcta es la \textbf{C}, ya que el error cometido por Marcos
fue no tener en cuenta que el segundo término tiene una variable
diferente, lo que resultó en la expresión \(-12a^2 + 8ab\).

Por lo tanto, la respuesta correcta es:

\[
\boxed{C}
\]

\begin{center}\rule{0.5\linewidth}{0.5pt}\end{center}

\subsection{Análisis de la Pregunta 2023-Matematicas-11-2-14, según
SABER 11
ICFES}\label{anuxe1lisis-de-la-pregunta-2023-matematicas-11-2-14-seguxfan-saber-11-icfes}

\subsection{Nivel de Desempeño: 3}\label{nivel-de-desempeuxf1o-3-2}

\subsection{Competencia:}\label{competencia-2}

\begin{itemize}
\tightlist
\item
  Argumentación

  \begin{itemize}
  \tightlist
  \item
    Aprendizaje: Valida procedimientos y estrategias matemáticas
    utilizadas para dar solución a problemas
  \item
    Evidencia: Argumenta a favor o en contra de un procedimiento para
    resolver un problema a la luz de criterios presentados o
    establecidos
  \end{itemize}
\end{itemize}

\subsection{Componente:
Numérico-Variacional}\label{componente-numuxe9rico-variacional-2}

\subsection{Estándar Asociado:}\label{estuxe1ndar-asociado-4}

Construyo expresiones algebraicas equivalentes a una expresión
algebraica dada.

\subsection{¿Qué evalúa?}\label{quuxe9-evaluxfaa-5}

La capacidad para identificar errores en procedimientos algebraicos de
factorización y operaciones con términos semejantes.

\subsection{Respuesta Correcta: C}\label{respuesta-correcta-c}

Justificación: El error está en el paso 3, donde Marcos no consideró que
al multiplicar \((−3a + 2b)4a\) el segundo término debe quedar con la
variable \(b\), resultando en \(-12a^2 + 8ab\) en lugar de
\(-12a^2 + 8a\).

\subsection{Distractores:}\label{distractores-3}

\begin{itemize}
\tightlist
\item
  A: Incorrecta porque la agrupación de términos semejantes en el paso 1
  está bien realizada
\item
  B: Incorrecta porque el cambio de signos en el paso 2 está bien
  aplicado
\item
  D: Incorrecta porque el error no está en la factorización final
\end{itemize}

\subsection{Contenidos Matemáticos
Curriculares:}\label{contenidos-matemuxe1ticos-curriculares-6}

\begin{itemize}
\tightlist
\item
  Álgebra y Cálculo

  \begin{itemize}
  \tightlist
  \item
    Período 2

    \begin{itemize}
    \tightlist
    \item
      Expresiones algebraicas
    \item
      Operaciones con expresiones algebraicas
    \end{itemize}
  \end{itemize}
\end{itemize}

\subsection{Tipo: No genérico}\label{tipo-no-genuxe9rico-1}

\subsection{Eje Axial Disciplinar:}\label{eje-axial-disciplinar-4}

Términos Semejantes, Factorización Básica: Álgebra

\subsection{Tarea:}\label{tarea-4}

Identificar el error en una secuencia de pasos algebraicos que involucra
operaciones con términos semejantes y multiplicación de expresiones
algebraicas.

\subsection{Grado: 8° - 9°}\label{grado-8---9}

\section{Pregunta 15}\label{pregunta-15}

\begin{center}\rule{0.5\linewidth}{0.5pt}\end{center}

\includegraphics{/home/manjaro_lenovo/Documentos/Math/2023-Matematicas-11-2/15/01.png}

\begin{center}\rule{0.5\linewidth}{0.5pt}\end{center}

\subsection{Análisis detallado del
problema}\label{anuxe1lisis-detallado-del-problema-6}

Identificamos los elementos clave del problema: - 2 opciones de entrada
- 4 opciones de plato fuerte - 3 opciones de postre - 4 opciones de
dulces de regalo

\subsection{Desarrollo Matemático}\label{desarrollo-matemuxe1tico}

\subsubsection{1. Estrategia de
Solución}\label{estrategia-de-soluciuxf3n}

Para encontrar el total de menús diferentes posibles, debemos aplicar el
principio multiplicativo de conteo, ya que: - Cada menú completo se
forma seleccionando una opción de cada categoría - Las selecciones son
independientes entre sí

\subsubsection{2. Aplicación del Principio
Multiplicativo}\label{aplicaciuxf3n-del-principio-multiplicativo}

\(\text{Total de menús} = \text{entradas} \times \text{platos fuertes} \times \text{postres} \times \text{dulces}\)

Sustituyendo los valores:

\(\text{Total} = 2 \times 4 \times 3 \times 4\)

\(\text{Total} = 96\)

\subsubsection{3. Verificación Lógica}\label{verificaciuxf3n-luxf3gica}

\begin{itemize}
\tightlist
\item
  Cada invitado debe elegir:

  \begin{itemize}
  \tightlist
  \item
    Una de las 2 entradas
  \item
    Uno de los 4 platos fuertes
  \item
    Uno de los 3 postres
  \item
    Uno de los 4 dulces de regalo
  \end{itemize}
\end{itemize}

\subsection{Conclusión}\label{conclusiuxf3n-11}

La respuesta correcta es C. 96 menús diferentes.

\subsubsection{Justificación:}\label{justificaciuxf3n-7}

El resultado de 96 menús diferentes se obtiene porque cada invitado
puede hacer su selección de manera independiente para cada categoría, y
el principio multiplicativo nos permite calcular todas las combinaciones
posibles mediante la multiplicación de las opciones disponibles en cada
categoría.

\subsection{Validación de Otras
Opciones:}\label{validaciuxf3n-de-otras-opciones}

\begin{itemize}
\tightlist
\item
  A. 24: Este número es muy pequeño y no considera todas las
  combinaciones posibles
\item
  B. 13: Este número es arbitrario y no refleja ninguna combinación
  lógica
\item
  D. 52: Este número es inferior al total real de combinaciones posibles
\end{itemize}

\begin{center}\rule{0.5\linewidth}{0.5pt}\end{center}

\subsection{Análisis de la Pregunta 2023-Matematicas-11-2-15, según
SABER 11
ICFES}\label{anuxe1lisis-de-la-pregunta-2023-matematicas-11-2-15-seguxfan-saber-11-icfes}

\begin{itemize}
\item
  Nivel de Desempeño: 3
\item
  Competencia: Formulación y Ejecución

  \begin{itemize}
  \tightlist
  \item
    Aprendizaje: Frente a un problema que involucre información
    cuantitativa, plantea e implementa estrategias que lleven a
    soluciones adecuadas.
  \item
    Evidencia: Ejecuta un plan de solución para un problema que
    involucra información cuantitativa o esquemática.
  \end{itemize}
\item
  Componente: Aleatorio
\item
  Estándar Asociado: Resuelvo y planteo problemas usando conceptos
  básicos de conteo y probabilidad (combinaciones, permutaciones,
  espacio muestral, muestreo aleatorio, muestreo con reemplazo).
\item
  ¿Qué evalúa?: La capacidad de aplicar el principio multiplicativo del
  conteo para determinar el número total de posibles combinaciones en
  una situación con múltiples opciones independientes.
\item
  Respuesta Correcta: C. 96\\
  Justificación: Al aplicar el principio multiplicativo del conteo, se
  multiplican las opciones disponibles de cada categoría:
  \(2 \times 4 \times 3 \times 4 = 96\) menús diferentes posibles.
\item
  Distractores:

  \begin{itemize}
  \tightlist
  \item
    A (24): Resulta de multiplicar incorrectamente solo algunas
    opciones.
  \item
    B (13): Número arbitrario sin fundamento matemático.
  \item
    D (52): Resultado de una operación incorrecta con los datos dados.
  \end{itemize}
\item
  Contenidos Matemáticos Curriculares:

  \begin{itemize}
  \tightlist
  \item
    Estadística
  \item
    Período 3

    \begin{itemize}
    \tightlist
    \item
      Conjuntos, Combinatoria y Probabilidad

      \begin{itemize}
      \tightlist
      \item
        Principio de multiplicación
      \end{itemize}
    \end{itemize}
  \end{itemize}
\item
  Genérico: No
\item
  Eje Axial Disciplinar: Eje 4 - Conjuntos, Casillas y Combinaciones
\item
  Tarea: Aplicar el principio multiplicativo para calcular el total de
  combinaciones posibles en una situación de selección múltiple
  independiente.
\item
  Grado: 11°
\end{itemize}

\section{Pregunta 16}\label{pregunta-16}

\begin{center}\rule{0.5\linewidth}{0.5pt}\end{center}

\includegraphics{/home/manjaro_lenovo/Documentos/Math/2023-Matematicas-11-2/16/01.png}

\begin{center}\rule{0.5\linewidth}{0.5pt}\end{center}

\subsection{Análisis detallado del
problema}\label{anuxe1lisis-detallado-del-problema-7}

En astronomía, para cuantificar el brillo de una estrella se utiliza la
magnitud absoluta \(M\) y la magnitud aparente \(m\). A partir de la
distancia \(d\) de la estrella, medida en pársecs, es posible relacionar
ambas magnitudes usando la siguiente ecuación:

\[ m - M = 5 \log_{10} \left( \frac{d}{10} \right) \]

Para una estrella con una distancia de 100 pársecs, la diferencia entre
la magnitud aparente y la magnitud absoluta es 5. Vamos a resolver el
problema paso a paso.

\subsubsection{Paso 1: Determinar la fórmula para la magnitud
aparente}\label{paso-1-determinar-la-fuxf3rmula-para-la-magnitud-aparente}

La fórmula dada es:

\[ m - M = 5 \log_{10} \left( \frac{d}{10} \right) \]

\subsubsection{Paso 2: Aplicar la fórmula con los datos
dados}\label{paso-2-aplicar-la-fuxf3rmula-con-los-datos-dados}

Queremos encontrar la magnitud aparente \(m\) de una estrella situada a
una distancia \(d\) de 1.000 pársecs cuya magnitud absoluta \(M\) es 15.

Sustituimos los valores en la fórmula:

\[ m - 15 = 5 \log_{10} \left( \frac{1000}{10} \right) \]

\subsubsection{Paso 3: Simplificar la expresión dentro del
logaritmo}\label{paso-3-simplificar-la-expresiuxf3n-dentro-del-logaritmo}

Simplificamos la fracción dentro del logaritmo:

\[ \frac{1000}{10} = 100 \]

Entonces la ecuación se convierte en:

\[ m - 15 = 5 \log_{10} (100) \]

\subsubsection{Paso 4: Calcular el
logaritmo}\label{paso-4-calcular-el-logaritmo}

Sabemos que:

\[ \log_{10} (100) = 2 \]

Sustituimos este valor en la ecuación:

\[ m - 15 = 5 \times 2 \]

\[ m - 15 = 10 \]

\subsubsection{\texorpdfstring{Paso 5: Resolver para
\(m\)}{Paso 5: Resolver para m}}\label{paso-5-resolver-para-m}

Sumamos 15 a ambos lados de la ecuación para encontrar \(m\):

\[ m = 10 + 15 \]

\[ m = 25 \]

\subsubsection{Conclusión}\label{conclusiuxf3n-12}

La magnitud aparente de la estrella es \textbf{25}. Por lo tanto, la
opción correcta es:

\textbf{A. 25}

\begin{center}\rule{0.5\linewidth}{0.5pt}\end{center}

\subsubsection{Análisis de la Pregunta 2023-Matematicas-11-2-16, según
SABER 11
ICFES}\label{anuxe1lisis-de-la-pregunta-2023-matematicas-11-2-16-seguxfan-saber-11-icfes}

\begin{itemize}
\item
  Nivel de Desempeño: 4
\item
  Competencia: Formulación y Ejecución

  \begin{itemize}
  \tightlist
  \item
    Aprendizaje: Frente a un problema que involucre información
    cuantitativa, plantea e implementa estrategias que lleven a
    soluciones adecuadas
  \item
    Evidencia: Ejecuta un plan de solución para un problema que
    involucra información cuantitativa o esquemática
  \end{itemize}
\item
  Componente: Numérico-Variacional
\item
  Estándar Asociado: Utilizo las técnicas de aproximación en procesos
  infinitos numéricos
\item
  ¿Qué evalúa?: La capacidad de aplicar una fórmula logarítmica en un
  contexto astronómico, realizando las sustituciones y operaciones
  algebraicas necesarias
\item
  Respuesta Correcta: A (25)
\item
  Justificación:\\
  Al sustituir los valores dados (\(d=1000\), \(M=15\))\\
  en la ecuación \(m - M = 5 \log_{10} \left(\frac{d}{10}\right)\),\\
  se obtiene \(m - 15 = 5 \log_{10}(100)\). ~

  Como \(\log_{10}(100)=2\),\\
  entonces \(m - 15 = 10\),\\
  de donde \(m = 25\)
\item
  Distractores:

  \begin{itemize}
  \tightlist
  \item
    B (30): Resultado de sumar incorrectamente 15 al valor 15 inicial
  \item
    C (0): Resultado de ignorar el término logarítmico
  \item
    D (5): Resultado de considerar solo el coeficiente 5 de la fórmula
  \end{itemize}
\item
  Contenidos Matemáticos Curriculares:

  \begin{itemize}
  \tightlist
  \item
    Álgebra y Cálculo
  \item
    Período 4

    \begin{itemize}
    \tightlist
    \item
      Funciones logarítmicas

      \begin{itemize}
      \tightlist
      \item
        Propiedades básicas
      \end{itemize}
    \end{itemize}
  \end{itemize}
\item
  No genérico
\item
  Eje Axial Disciplinar: Funciones II (Comportamientos Funcionales)
\item
  Tarea: Resolver una ecuación logarítmica aplicada al contexto
  astronómico mediante sustitución y operaciones algebraicas
\item
  Grado: 11°
\end{itemize}

\section{Pregunta 17}\label{pregunta-17}

\begin{center}\rule{0.5\linewidth}{0.5pt}\end{center}

\includegraphics{/home/manjaro_lenovo/Documentos/Math/2023-Matematicas-11-2/17/01.png}

\begin{center}\rule{0.5\linewidth}{0.5pt}\end{center}

\subsection{Análisis detallado del
problema}\label{anuxe1lisis-detallado-del-problema-8}

La opción correcta es la \textbf{D. \(n^2\)}.

\textbf{Paso a paso:}

\begin{enumerate}
\def\labelenumi{\arabic{enumi}.}
\item
  \textbf{Identificar la secuencia de puntos:} El enunciado describe una
  secuencia de puntos ganados con cada carta verde:

  \begin{itemize}
  \tightlist
  \item
    Primera carta: 1 punto
  \item
    Segunda carta: 3 puntos
  \item
    Tercera carta: 5 puntos
  \item
    \ldots{} y así sucesivamente.
  \end{itemize}

  Podemos observar que los puntos ganados forman una secuencia de
  números impares.
\item
  \textbf{Encontrar una fórmula para los puntos de la \emph{n}-ésima
  carta:} Analizando la secuencia, podemos deducir una fórmula para la
  cantidad de puntos ganados con la \emph{n}-ésima carta verde.
  Observemos la relación:

  \begin{itemize}
  \tightlist
  \item
    Para la 1ª carta (n=1): 1 punto (2*1 - 1)
  \item
    Para la 2ª carta (n=2): 3 puntos (2*2 - 1)
  \item
    Para la 3ª carta (n=3): 5 puntos (2*3 - 1)
  \item
    Para la \emph{n}-ésima carta: \(2n - 1\) puntos
  \end{itemize}
\item
  \textbf{Calcular el puntaje total:} El puntaje total se obtiene
  sumando los puntos de todas las cartas verdes obtenidas. Si el jugador
  obtuvo \emph{n} cartas verdes, necesitamos sumar los puntos de la
  primera carta hasta la \emph{n}-ésima carta. Esto sería la suma de la
  siguiente serie:

  \(1 + 3 + 5 + ... + (2n - 1)\)

  Esta es la suma de los primeros \emph{n} números impares.
\item
  \textbf{Aplicar la fórmula para la suma de los primeros \emph{n}
  números impares:} Existe una fórmula matemática que establece que la
  suma de los primeros \emph{n} números impares es igual a \(n^2\).

  \emph{Demostración (opcional):} Podemos ver esto como una serie
  aritmética donde el primer término es \(a_1 = 1\), la diferencia común
  es \(d = 2\), y el número de términos es \(n\). La suma de una serie
  aritmética se calcula con la fórmula:

  \(S_n = \frac{n}{2} [2a_1 + (n-1)d]\)

  Sustituyendo los valores:

  \(S_n = \frac{n}{2} [2(1) + (n-1)2]\)
  \(S_n = \frac{n}{2} [2 + 2n - 2]\) \(S_n = \frac{n}{2} [2n]\)
  \(S_n = n^2\)
\item
  \textbf{Verificar las otras opciones:}

  \begin{itemize}
  \tightlist
  \item
    \textbf{A. \(2n - 1\):} Esta fórmula representa los puntos ganados
    con la \emph{n}-ésima carta, no el puntaje total.
  \item
    \textbf{B. \(2n + 1\):} Esta fórmula no coincide con la secuencia de
    puntos ganados por las cartas.
  \item
    \textbf{C. \(n + 2\):} Esta fórmula tampoco coincide con la
    secuencia de puntos ganados ni con el puntaje total.
  \end{itemize}
\end{enumerate}

\textbf{Conclusión:} La fórmula que permite calcular el puntaje total de
un jugador que obtuvo \emph{n} cartas verdes es \textbf{\(n^2\)}. Cada
carta sucesiva aporta un número impar de puntos, y la suma de los
primeros \emph{n} números impares es igual a \(n^2\).

\begin{center}\rule{0.5\linewidth}{0.5pt}\end{center}

\subsubsection{Análisis de la Pregunta 2023-Matematicas-11-2-17, según
SABER 11
ICFES}\label{anuxe1lisis-de-la-pregunta-2023-matematicas-11-2-17-seguxfan-saber-11-icfes}

\begin{itemize}
\item
  Nivel de Desempeño: 4
\item
  Competencia: Interpretación y Representación

  \begin{itemize}
  \tightlist
  \item
    Afirmación: Comprende y transforma la información cuantitativa y
    esquemática presentada en distintos formatos
  \item
    Evidencia: Da cuenta de las características básicas de la
    información presentada en diferentes formatos
  \end{itemize}
\item
  Componente: Numérico-Variacional
\item
  Estándar Asociado: Reconozco y generalizo propiedades de las
  relaciones entre números racionales y de las operaciones entre ellos
  en diferentes contextos
\item
  ¿Qué evalúa?: La capacidad para identificar patrones numéricos y
  traducirlos a una expresión algebraica que generaliza la suma de una
  secuencia aritmética de números impares
\item
  Respuesta Correcta: D. \(n^2\)
\end{itemize}

Justificación:

La secuencia de puntos forma una progresión aritmética de números
impares: 1, 3, 5, 7, \ldots, \((2n-1)\)

La suma de los primeros \(n\) números impares viene dada por la fórmula
\(n^2\), lo que se puede verificar:

\begin{itemize}
\tightlist
\item
  Para n=1: \(1 = 1^2\)
\item
  Para n=2: \(1 + 3 = 2^2\)
\item
  Para n=3: \(1 + 3 + 5 = 3^2\)
\item
  Para n=4: \(1 + 3 + 5 + 7 = 4^2\)
\end{itemize}

Motivos de las opciones incorrectas:\\
- A. \(2n-1\): Representa el último término de la secuencia, no la
suma\\
- B. \(2n+1\): Fórmula incorrecta que no corresponde al patrón\\
- C. \(n+2\): No representa el comportamiento cuadrático de la suma

\begin{itemize}
\tightlist
\item
  Contenidos Matemáticos Curriculares:

  \begin{itemize}
  \tightlist
  \item
    Álgebra y Cálculo
  \item
    Período 2

    \begin{itemize}
    \tightlist
    \item
      Expresiones algebraicas

      \begin{itemize}
      \tightlist
      \item
        Operaciones con expresiones algebraicas
      \end{itemize}
    \end{itemize}
  \end{itemize}
\item
  No genérico
\item
  Eje Axial Disciplinar: Eje 3 - Función Lineal, Generalizaciones y
  Grado 2
\item
  Tarea: Identificar la expresión algebraica que generaliza la suma de
  una secuencia aritmética de números impares
\item
  Grado: 11°
\end{itemize}

\section{Pregunta 18}\label{pregunta-18}

\begin{center}\rule{0.5\linewidth}{0.5pt}\end{center}

\includegraphics{/home/manjaro_lenovo/Documentos/Math/2023-Matematicas-11-2/18/01.png}

\begin{center}\rule{0.5\linewidth}{0.5pt}\end{center}

\subsection{Análisis detallado del
problema}\label{anuxe1lisis-detallado-del-problema-9}

Para resolver este problema, es importante entender cómo se calcula la
masa de una sustancia a partir de su volumen y densidad. La fórmula para
calcular la masa es:

\[ \text{masa} = \text{densidad} \times \text{volumen} \]

Dado que Jaime dispone de 0,01 litros de mercurio y la densidad del
mercurio es de \(13.600 \, \text{kg/m}^3\), necesitamos convertir el
volumen a metros cúbicos (\(\text{m}^3\)) para que las unidades
coincidan.

Vamos a analizar cada opción paso a paso:

\subsubsection{Opción A}\label{opciuxf3n-a}

\textbf{No, puesto que se equivocó al multiplicar decimales.}

Esta opción sugiere que Jaime se equivocó en la multiplicación de
decimales. Sin embargo, el error principal no está en la multiplicación
de decimales, sino en la conversión de unidades.

\subsubsection{Opción B}\label{opciuxf3n-b}

\textbf{Sí, puesto que 0,01 litros equivale a un metro cúbico.}

Esta opción es incorrecta. 0,01 litros no equivale a un metro cúbico. La
conversión correcta es:

\[ 1 \, \text{litro} = 0,001 \, \text{m}^3 \]

Por lo tanto, 0,01 litros es:

\[ 0,01 \, \text{litros} = 0,01 \times 0,001 \, \text{m}^3 = 0,00001 \, \text{m}^3 \]

\subsubsection{Opción C}\label{opciuxf3n-c}

\textbf{No, puesto que le faltó hacer la conversión de litros a metros
cúbicos.}

Esta es la opción correcta. Jaime no realizó la conversión de litros a
metros cúbicos, lo cual es necesario para que las unidades coincidan en
la fórmula de la masa.

\subsubsection{Opción D}\label{opciuxf3n-d}

\textbf{Sí, puesto que, para hallar la masa, basta con multiplicar la
densidad por el volumen.}

Esta opción es incorrecta porque, aunque la fórmula para hallar la masa
es correcta, Jaime no realizó la conversión de unidades necesaria.

\subsubsection{Resolución correcta}\label{resoluciuxf3n-correcta}

\begin{enumerate}
\def\labelenumi{\arabic{enumi}.}
\tightlist
\item
  \textbf{Convertir el volumen de litros a metros cúbicos:}
\end{enumerate}

\[ 0,01 \, \text{litros} = 0,01 \times 0,001 \, \text{m}^3 = 0,00001 \, \text{m}^3 \]

\begin{enumerate}
\def\labelenumi{\arabic{enumi}.}
\setcounter{enumi}{1}
\tightlist
\item
  \textbf{Calcular la masa usando la densidad y el volumen convertido:}
\end{enumerate}

\[ \text{masa} = \text{densidad} \times \text{volumen} \]
\[ \text{masa} = 13.600 \, \text{kg/m}^3 \times 0,00001 \, \text{m}^3 \]
\[ \text{masa} = 0,136 \, \text{kg} \]

Por lo tanto, el peso correcto del mercurio es 0,136 kg, no 136 kg.

\subsubsection{Conclusión}\label{conclusiuxf3n-13}

La opción correcta es la \textbf{C}. `Jaime se equivocó porque no
realizó la conversión de litros a metros cúbicos'.

\begin{center}\rule{0.5\linewidth}{0.5pt}\end{center}

\subsubsection{Análisis de la Pregunta 2023-Matematicas-11-2-18, según
SABER 11
ICFES}\label{anuxe1lisis-de-la-pregunta-2023-matematicas-11-2-18-seguxfan-saber-11-icfes}

\begin{itemize}
\tightlist
\item
  Nivel de Desempeño: 3
\item
  Competencia: Argumentación

  \begin{itemize}
  \tightlist
  \item
    Aprendizaje: Valida procedimientos y estrategias matemáticas
    utilizadas para dar solución a problemas
  \item
    Evidencia: Argumenta a favor o en contra de un procedimiento para
    resolver un problema a la luz de criterios presentados o
    establecidos
  \end{itemize}
\item
  Componente: Numérico-Variacional
\item
  Estándar Asociado: Justifico la pertinencia de utilizar unidades de
  medida estandarizadas en situaciones tomadas de distintas ciencias.
\item
  ¿Qué evalúa?: La capacidad para identificar errores en procedimientos
  que involucran conversión de unidades de volumen y cálculos con
  densidad.
\item
  Respuesta Correcta: C
\item
  Justificación: El error está en no convertir 0,01 litros a metros
  cúbicos antes de multiplicar por la densidad. La conversión correcta
  sería:
\end{itemize}

\(0,01 \text{ L} = 0,01 \times 0,001 \text{ m}^3 = 0,00001 \text{ m}^3\)

Luego: \(\text{masa} = \text{densidad} \times \text{volumen}\)
\(\text{masa} = 13.600 \text{ kg/m}^3 \times 0,00001 \text{ m}^3 = 0,136 \text{ kg}\)

\begin{itemize}
\tightlist
\item
  Distractores:

  \begin{itemize}
  \tightlist
  \item
    A: Asume que el error está en la multiplicación de decimales
  \item
    B: Asume incorrectamente la equivalencia entre litros y metros
    cúbicos
  \item
    D: Aunque la fórmula es correcta, ignora la necesidad de conversión
    de unidades
  \end{itemize}
\item
  Contenidos Matemáticos Curriculares:

  \begin{itemize}
  \tightlist
  \item
    Álgebra y Cálculo
  \item
    Período 1

    \begin{itemize}
    \tightlist
    \item
      Números y Operaciones

      \begin{itemize}
      \tightlist
      \item
        Uso de números y operaciones en la resolución de problemas
      \end{itemize}
    \end{itemize}
  \end{itemize}
\item
  Genérico: Sí
\item
  Eje Axial Disciplinar: Eje 1 - Paso 1, Paso 2: Seguimiento de
  instrucciones
\item
  Tarea: Validar la pertinencia de un procedimiento de cálculo que
  involucra conversión de unidades y densidad
\item
  Grado: 9° - 11°
\end{itemize}

\section{Pregunta 19}\label{pregunta-19}

\begin{center}\rule{0.5\linewidth}{0.5pt}\end{center}

\includegraphics{/home/manjaro_lenovo/Documentos/Math/2023-Matematicas-11-2/19/all.png}

\begin{center}\rule{0.5\linewidth}{0.5pt}\end{center}

\subsection{Análisis detallado del
problema}\label{anuxe1lisis-detallado-del-problema-10}

Pablo quiere ir desde el punto \(M\) hasta el \(N\) y tiene dos caminos
posibles como se muestra en la figura. Pablo afirmó que la distancia
entre \(M\) y \(N\) es igual a 60 metros por cualquiera de los dos
caminos. Vamos a verificar esta afirmación paso a paso.

\subsection{Análisis de las
Opciones}\label{anuxe1lisis-de-las-opciones-2}

\subsubsection{Opción A}\label{opciuxf3n-a-1}

\textbf{Afirmación:} Sí, porque entre ambos recorridos se encierran seis
cuadrados de 10 metros de lado.

\textbf{Análisis:} Esta opción sugiere que la distancia total es 60
metros porque se encierran seis cuadrados de 10 metros de lado. Sin
embargo, esto no considera el recorrido específico de cada camino.

\subsubsection{Opción B}\label{opciuxf3n-b-1}

\textbf{Afirmación:} No, porque se debe multiplicar el perímetro de cada
cuadrado por los seis segmentos recorridos para conocer el recorrido
total.

\textbf{Análisis:} Esta opción sugiere que se debe considerar el
perímetro de cada cuadrado multiplicado por los seis segmentos. Esto no
es correcto porque el perímetro de un cuadrado no es relevante para
calcular la distancia total recorrida.

\subsubsection{Opción C}\label{opciuxf3n-c-1}

\textbf{Afirmación:} Sí, porque cada lado de un cuadrado es igual a 10
metros, y se deben recorrer seis lados para ir desde \(M\) hasta \(N\)
en ambas rutas.

\textbf{Análisis:} Esta opción es correcta. Vamos a verificarlo:

\begin{enumerate}
\def\labelenumi{\arabic{enumi}.}
\item
  \textbf{Camino 1 (línea discontinua):}

  \begin{itemize}
  \tightlist
  \item
    Desde \(M\) hacia la derecha: 10 metros.
  \item
    Hacia abajo: 10 metros.
  \item
    Hacia la derecha: 10 metros.
  \item
    Hacia abajo: 10 metros.
  \item
    Hacia la derecha: 10 metros.
  \item
    Hacia abajo: 10 metros.
  \end{itemize}

  Total: \(10 + 10 + 10 + 10 + 10 + 10 = 60\) metros.
\item
  \textbf{Camino 2 (línea punteada):}

  \begin{itemize}
  \tightlist
  \item
    Desde \(M\) hacia abajo: 10 metros.
  \item
    Hacia la derecha: 10 metros.
  \item
    Hacia abajo: 10 metros.
  \item
    Hacia la derecha: 10 metros.
  \item
    Hacia abajo: 10 metros.
  \item
    Hacia la derecha: 10 metros.
  \end{itemize}

  Total: \(10 + 10 + 10 + 10 + 10 + 10 = 60\) metros.
\end{enumerate}

\subsubsection{Opción D}\label{opciuxf3n-d-1}

\textbf{Afirmación:} No, porque se debe dividir el perímetro de un
cuadrado entre la medida de un lado, y este resultado multiplicarlo por
los seis segmentos.

\textbf{Análisis:} Esta opción sugiere dividir el perímetro de un
cuadrado entre la medida de un lado y luego multiplicar por los seis
segmentos. Esto no es correcto porque el perímetro de un cuadrado no es
relevante para calcular la distancia total recorrida.

\subsection{Conclusión}\label{conclusiuxf3n-14}

La opción correcta es la \textbf{C}. Cada lado de un cuadrado es igual a
10 metros, y se deben recorrer seis lados para ir desde \(M\) hasta
\(N\) en ambas rutas, resultando en una distancia total de 60 metros
para cualquiera de los dos caminos.

\[ \boxed{C} \]

\begin{center}\rule{0.5\linewidth}{0.5pt}\end{center}

\subsubsection{Análisis de la Pregunta 2023-Matematicas-11-2-19, según
SABER 11
ICFES}\label{anuxe1lisis-de-la-pregunta-2023-matematicas-11-2-19-seguxfan-saber-11-icfes}

\begin{itemize}
\tightlist
\item
  Nivel de Desempeño: 3
\item
  Competencia: Argumentación

  \begin{itemize}
  \tightlist
  \item
    Afirmación: Valida procedimientos y estrategias matemáticas
    utilizadas para dar solución a problemas
  \item
    Evidencia: Plantea afirmaciones que sustentan o refutan una
    interpretación dada a la información disponible
  \end{itemize}
\item
  Componente: Geométrico-Métrico
\item
  Estándar Asociado: Generalizo procedimientos de cálculo válidos para
  encontrar el área de regiones planas y el volumen de sólidos.
\item
  ¿Qué evalúa?: La capacidad para validar una afirmación sobre
  distancias en una cuadrícula usando propiedades geométricas básicas.
\item
  Respuesta Correcta: C, porque cada lado del cuadrado mide 10 metros y
  en ambas rutas se deben recorrer exactamente 6 lados para ir de M a N,
  resultando en 60 metros totales.
\end{itemize}

Distractores:

\begin{itemize}
\tightlist
\item
  A: Confunde el área encerrada con la distancia recorrida
\item
  B: Aplica incorrectamente el concepto de perímetro
\item
  D: Realiza operaciones sin sentido geométrico
\end{itemize}

Contenidos Matemáticos Curriculares:

\begin{itemize}
\tightlist
\item
  Geometría
\item
  Período 1

  \begin{itemize}
  \tightlist
  \item
    Geometría Básica

    \begin{itemize}
    \tightlist
    \item
      Medición de distancias
    \item
      Trayectorias en el plano cartesiano
    \end{itemize}
  \end{itemize}
\end{itemize}

Tipo: No genérico

Eje Axial Disciplinar: Complementarios, Suplementarios y Operaciones
Proporcionales

Tarea: Validar una afirmación sobre la igualdad de distancias entre dos
puntos siguiendo diferentes trayectorias rectilíneas en una cuadrícula.

Grado sugerido: 9° - 11°

\section{Pregunta 20}\label{pregunta-20}

\begin{center}\rule{0.5\linewidth}{0.5pt}\end{center}

\includegraphics{/home/manjaro_lenovo/Documentos/Math/2023-Matematicas-11-2/20/01.png}

\begin{center}\rule{0.5\linewidth}{0.5pt}\end{center}

\subsection{Análisis detallado del
problema}\label{anuxe1lisis-detallado-del-problema-11}

Tenemos dos polinomios cuadráticos:

\begin{itemize}
\tightlist
\item
  \(P_1(x) = x^2 - 2x + 3\)
\item
  \(P_2(x) = x^2 + 4x + 1\)
\end{itemize}

Daniela afirma que el máximo exponente en el producto será 4, sin
necesidad de realizar la multiplicación.

\subsection{Análisis Paso a Paso}\label{anuxe1lisis-paso-a-paso}

\begin{enumerate}
\def\labelenumi{\arabic{enumi})}
\tightlist
\item
  Al multiplicar dos polinomios, los exponentes de los términos se
  suman. Por ejemplo:

  \begin{itemize}
  \tightlist
  \item
    \((x^a)(x^b) = x^{a+b}\)
  \end{itemize}
\item
  En este caso, el máximo exponente en cada polinomio es 2, entonces:

  \begin{itemize}
  \tightlist
  \item
    Máximo exponente de \(P_1(x)\) = 2
  \item
    Máximo exponente de \(P_2(x)\) = 2
  \end{itemize}
\item
  Al multiplicar los términos con los exponentes más altos:

  \begin{itemize}
  \tightlist
  \item
    \((x^2)(x^2) = x^{2+2} = x^4\)
  \end{itemize}
\item
  Por lo tanto, el máximo exponente en el producto será:

  \begin{itemize}
  \tightlist
  \item
    2 + 2 = 4
  \end{itemize}
\end{enumerate}

\subsection{Conclusión}\label{conclusiuxf3n-15}

La respuesta correcta es la opción \textbf{A}: ``Porque el exponente 4
corresponde a la suma de los máximos exponentes en cada polinomio.''

\subsection{Justificación de por qué las otras opciones son
incorrectas:}\label{justificaciuxf3n-de-por-quuxe9-las-otras-opciones-son-incorrectas}

\begin{itemize}
\tightlist
\item
  \textbf{B} es incorrecta porque sugiere multiplicar los exponentes (2
  × 2 = 4), lo cual no es la regla correcta.
\item
  \textbf{C} es incorrecta porque el número de términos no determina el
  exponente máximo.
\item
  \textbf{D} es incorrecta porque no se eleva al cuadrado el exponente
  máximo.
\end{itemize}

\subsection{Principio Matemático
General}\label{principio-matemuxe1tico-general}

En el álgebra de polinomios, cuando multiplicamos dos términos:\\
\(x^m \cdot x^n = x^{m+n}\)

Este principio fundamental explica por qué el máximo exponente en el
producto será la suma de los exponentes máximos de los polinomios
factores.

\begin{center}\rule{0.5\linewidth}{0.5pt}\end{center}

\subsubsection{Análisis de la Pregunta 2023-Matematicas-11-2-20, según
SABER 11
ICFES}\label{anuxe1lisis-de-la-pregunta-2023-matematicas-11-2-20-seguxfan-saber-11-icfes}

\begin{itemize}
\tightlist
\item
  Nivel de Desempeño: 3
\item
  Competencia: Argumentación

  \begin{itemize}
  \tightlist
  \item
    Aprendizaje: Valida procedimientos y estrategias matemáticas
    utilizadas para dar solución a problemas
  \item
    Evidencia: Plantea afirmaciones que sustentan o refutan una
    interpretación dada a la información disponible
  \end{itemize}
\item
  Componente: Numérico-Variacional
\item
  Estándar Asociado: Utilizo las propiedades de los números reales para
  justificar procedimientos y diferentes representaciones de
  subconjuntos de ellos
\item
  ¿Qué evalúa?: La comprensión de las propiedades de los exponentes al
  multiplicar polinomios
\end{itemize}

\subsection{Respuesta Correcta y
Justificación}\label{respuesta-correcta-y-justificaciuxf3n-1}

La opción correcta es A: ``Porque el exponente 4 corresponde a la suma
de los máximos exponentes en cada polinomio''

\subsubsection{Justificación}\label{justificaciuxf3n-8}

Al multiplicar dos polinomios \(P_1(x) = x^2 - 2x + 3\) y
\(P_2(x) = x^2 + 4x + 1\):

\begin{enumerate}
\def\labelenumi{\arabic{enumi})}
\tightlist
\item
  Los exponentes se suman al multiplicar términos:
  \(x^a \cdot x^b = x^{a+b}\)
\item
  El mayor exponente en cada polinomio es 2
\item
  Por tanto: \(x^2 \cdot x^2 = x^{2+2} = x^4\)
\end{enumerate}

\subsubsection{Distractores}\label{distractores-4}

B: Incorrecta - Confunde suma con producto de exponentes C: Incorrecta -
Relaciona erróneamente cantidad de términos con exponente máximo D:
Incorrecta - Malinterpreta las reglas de exponentes

\subsection{Contenidos Matemáticos
Curriculares}\label{contenidos-matemuxe1ticos-curriculares-7}

\begin{itemize}
\tightlist
\item
  Álgebra y Cálculo
\item
  Período: 2

  \begin{itemize}
  \tightlist
  \item
    Álgebra

    \begin{itemize}
    \tightlist
    \item
      Operaciones con expresiones algebraicas
    \end{itemize}
  \end{itemize}
\item
  No genérico
\item
  Eje Axial Disciplinar: Términos Semejantes, Factorización Básica:
  Álgebra
\item
  Tarea: Justificar el exponente máximo resultante al multiplicar dos
  polinomios usando las propiedades de los exponentes
\item
  Grado: 8° - 9°
\end{itemize}

\section{Referencias}\label{referencias}

\section{Referencias}\label{referencias-1}

\phantomsection\label{refs}
\begin{CSLReferences}{0}{1}
\end{CSLReferences}

\end{document}
